% refined Kauffman 不変量の large N 双対（日本語の準備稿）
% 構成：docs/refined_kauffman/paper/outline.md、主張の格付け：claims.md（主張の直前に % [ID] を書く）
% コンパイル：cd docs/refined_kauffman/paper/ja && TEXMFVAR=$TMPDIR/texmfvar lualatex main.tex && bibtex main && lualatex main.tex（2 回）
% 参考文献：refs.bib（biblatex、alphabetic 様式：引用は [AS12] のように著者の略号＋年。Young 図の [1] と区別するため）
\documentclass[11pt,a4paper]{ltjsarticle}
\usepackage{amsmath,amssymb,amsthm,bm}
\usepackage{booktabs,array}
\usepackage[margin=25mm]{geometry}
\usepackage[backend=bibtex,style=alphabetic,maxalphanames=4,minalphanames=4,doi=true,eprint=true]{biblatex}
\addbibresource{refs.bib}
\usepackage[hidelinks]{hyperref}
\allowdisplaybreaks
\sloppy

\newtheorem{theorem}{定理}[section]
\newtheorem{proposition}[theorem]{命題}
\newtheorem{lemma}[theorem]{補題}
\newtheorem{corollary}[theorem]{系}
\newtheorem{conjecture}[theorem]{予想}
\theoremstyle{definition}
\newtheorem{definition}[theorem]{定義}
\newtheorem{hypothesis}[theorem]{仮説}
\newtheorem{compresult}[theorem]{計算結果}
\theoremstyle{remark}
\newtheorem{remark}[theorem]{注意}

\newcommand{\ba}{\bm{a}}
\newcommand{\bq}{\bm{q}}
\newcommand{\bt}{\bm{t}}
\newcommand{\tlam}{\tilde\lambda}
\newcommand{\ql}{q_\ell}
\newcommand{\qr}{q_r}
\newcommand{\Rred}{\mathcal{R}}
\newcommand{\SO}{\mathrm{SO}}
\newcommand{\U}{\mathrm{U}}
\newcommand{\cov}{\mathrm{cov}}
\newcommand{\open}{\textbf{【未解決】}}

\title{refined Kauffman 不変量の large $N$ 双対：\\orientifold の refined 指数とクロスキャップの構造\\
\large（日本語の準備稿）}
\author{亀山 昌也（Masaya Kameyama）}
\date{2026年10月20日（準備稿）}

\begin{document}
\maketitle
\begin{abstract}
SO$(N)$ Chern--Simons 理論の彩色 Kauffman 不変量の large $N$ 双対（Mariño）を、Aganagic--Schaeffer の orientifold の refined 位相的弦の枠組みで refined にする。
SO$(2N)$ の refined Chern--Simons 不変量から、双対 Cauchy 核（t-brane）と、行列式と荷電共役で不変な規格化を用いて、refined のクロスキャップ部分 $g = f - \tfrac12 f^{\cov}$ を取り出す。

主な結果は次の 5 つである。
\begin{enumerate}
\item \textbf{構造}：M 理論の指数で、orientifold の一粒子状態を $\tilde I$ で入れ替わる対・$\tilde I$ で不変な状態・凍結した状態に分けると、全 $d$ の構造 $g = \tfrac12 E/(1+t) + C$ が従う。$E$ と $C$ の分離は指数からは決まらない。閉弦では、SO$(2N)$ の分配関数の積表示から作った自由エネルギーが、$N$ の 1 次式を除いて Aganagic--Schaeffer の式と一致することを確かめた。また、Avetisyan--Mkrtchyan の「Klein 瓶」の寄与が、この分解の並べ替えであることを示す。
\item \textbf{計算}：DAHA-Jones 多項式を一般の $t$ で直接計算し、refined SO$(2N)$ の trefoil の 2 箱・3 箱と $T(2,5)$ の 2 箱を一般の $N$ で求める。大きな $N$ のクロスキャップ部分は整数係数の Laurent 多項式で、$(1+t)$ の極は Adams の予言どおり $p_2$ 方向にだけ現れる。$t = q$ では Mariño の結果を全係数で再現する。
\item \textbf{普遍因子}：refined だけの部分 $\hat g - \hat g|_{t=q}(\lambda \to \lambda (q/t)^{1/2})$ は、[1]・2 箱・3 箱・Hopf 絡み目のすべてで $(\bt^2 - 1)(\ba^2 - \bq^2)$ で割り切れる。Kähler 変数のずれの向きは一意に決まる。
\item \textbf{正値性}：トーラス結び目の [1] では、$G$ 展開の後に符号がそろう判定式が成り立つ。多箱では符号がそろわない。この障害を、凍結した状態の多重項の基底と $E$・$C$ の分離として定式化する。$E$ と $C$ の分離が指数から決まらないことを示し、composite の DAHA の値からも状態の数が得られなかったことを報告する。多箱の $c = 1$ 部分が $\sigma$ 対称でない理由を、unrefined の段階での $c$ の取り違え（hatted では $c = 1$ は $\lambda^\beta$ の $\beta \equiv |R| - 1 \pmod 2$ の部分）と、refined に固有の $\sigma$ 奇部分（$(\bt^2-1)(\ql - \ql^{-1})$ で割り切れる）に分ける。
\item \textbf{Sp$(2N)$}：C$_N$ 型の DAHA を実装し、Avetisyan--Mkrtchyan の refined Sp$(2N)$（$u = t^2$）の結び目不変量を初めて計算した。[1] では $g^{\mathrm{Sp}}(\ba) = -g^{\SO}(-\ba)$（$c = 1$ の符号反転が refined でも成り立つ）。2 箱では彩色の $W_R$ の段階で $W^{\mathrm{Sp}}_R(\ba) = (-1)^{|R|}W^{\SO}_R(-\ba)$ が成り立ち、クロスキャップ部分の差は $p_2$ の定数項と Adams 項に由来する閉じた形 $\pm\psi^2[\delta]$ に尽きる。
\end{enumerate}
\end{abstract}

\tableofcontents

\section{序論}\label{sec:intro}

\subsection{背景}
U$(N)$ Chern--Simons 理論の彩色 HOMFLY-PT 不変量は、resolved conifold 上の開弦の位相的弦と large $N$ で双対であり、Ooguri--Vafa の再定式化によって整数の BPS 不変量（LMOV 不変量）で書ける。
SO$(N)$・Sp$(N)$ の場合は、Calabi--Yau の orientifold の位相的弦が対応する。Mariño \cite{Marino09} は、SO$(N)$ の彩色 Kauffman 不変量から、U$(N)$ の composite 表現の寄与の半分を引いた残りが、クロスキャップ $c = 1, 2$ の整数 BPS 不変量で書けることを予想し、多くの例で確かめた（Bouchard--Florea--Mariño \cite{BFM04,BFM05}、Nawata--Ramadevi--Zodinmawia \cite{NRZ13} なども参照）。

refined の側では、Aganagic--Shakirov \cite{AgSh11} が refined Chern--Simons 理論を定式化した。著者と Nawata \cite{KN17} は、U$(N)$ の refined Chern--Simons 不変量が resolved conifold 上の refined BPS 不変量を与え、トーラス結び目で正の整数になることを確かめた（refined large $N$ 双対）。
Aganagic--Schaeffer \cite{AgSc12} は、orientifold の refined 位相的弦を M 理論の指数として定義した。そのうえで SO$(2N)$ の refined Chern--Simons 不変量を計算し、大きな $N$ の閉弦の自由エネルギーが第二量子化の整数性を満たすことを示した。
ただし開弦（結び目）については、refined の BPS 不変量を取り出す式も、その整数性・正値性の検証も無かった。
著者は以前（2018 年）に、Mariño の構造を素朴に refined にする計算を行ったが、整数性が得られなかった。

\subsection{本論文の問い}
本論文では次の問いを扱う。
\begin{enumerate}
\item SO$(2N)$ の refined Chern--Simons 不変量から、refined のクロスキャップの BPS 不変量を取り出す式はどう書けるか。
\item その整数性はどの範囲で成り立つか。
\item refined で正値性は見えるか。見えないなら、何が障害か。
\end{enumerate}

\subsection{結果の概要}
\begin{itemize}
\item \textbf{定式化（§\ref{sec:prelim}）}：Mariño 型の bosonic な核を素朴に refined にすると、unknot の 2 箱で既に $U = t^N$ の極が打ち消さない。代わりに双対 Cauchy 核（t-brane）を使うと、unknot の積表示が自動的に成り立つ。
規格化は、U$(n)$ の不変量が行列式と荷電共役で不変になる「一様な規格化」をとる。reformulation には Mariño の $M(q^{-1})$ を使う。
\item \textbf{構造（§\ref{sec:index}）}：orientifold の M 理論の指数から、全 $d$ の構造 $g = \tfrac12 E/(1+t) + C$ を導く（命題 \ref{prop:alld}）。$E$ と $C$ の分離は指数からは決まらない（命題 \ref{prop:EC}）。
この構造は、Witten \cite{Witten11} の Chern--Simons 分配関数の trace 表示の SO$(2N)$ 版として読める（命題 \ref{prop:trace}）。境界ゼロモードへの $\tilde I$ の作用から、$E$ の境界因子が $z_q^{n-1}(z_q^+)^n$ になることを導き、Hopf 絡み目・結び目・framed unknot で確かめた（命題 \ref{prop:boundary}）。
閉弦では、SO$(2N)$ の分配関数の積表示が $N$ の 1 次式を除いて Aganagic--Schaeffer の式と一致することを確かめ、Avetisyan--Mkrtchyan \cite{AM22} の「Klein 瓶」の寄与がこの分解の並べ替えであることを示す（命題 \ref{prop:AM}）。
\item \textbf{unknot と Hopf 絡み目（§\ref{sec:hopf}）}：整数性が成り立つ。Kähler 変数は $\lambda \to \tlam = \lambda(q/t)^{1/2}$ とずれ、$E_{22}, E_{24}$ は一意に決まる。
\item \textbf{トーラス結び目の [1]（§\ref{sec:torus}）}：$T(2,2m+1)$（$m \le 5$）で、$c = 1$ 部分の $G$ 展開と $c = 2$ 部分の閉じた形が得られ、符号がそろう判定式が成り立つ。
\item \textbf{一般の $N$ の多箱（§\ref{sec:generalN}）}：DAHA-Jones 多項式を一般の $t$ で計算し、refined SO$(2N)$ の trefoil の 2 箱・3 箱と $T(2,5)$ の 2 箱を一般の $N$ で求めた。
クロスキャップ部分は整数で、$(1+t)$ の極は $p_2$ 方向にだけ現れる。refined だけの部分は普遍因子 $(\bt^2 - 1)(\ba^2 - \bq^2)$ で割り切れる。
\item \textbf{refined Sp$(2N)$（§\ref{sec:sp}）}：C$_N$ 型の DAHA を実装し、Avetisyan--Mkrtchyan の refined Sp$(2N)$ の trefoil・$T(2,5)$ の [1] と trefoil の 2 箱を一般の $N$ で求めた。[1] では $g^{\mathrm{Sp}}(\ba) = -g^{\SO}(-\ba)$、2 箱では $W^{\mathrm{Sp}}_R(\ba) = (-1)^{|R|} W^{\SO}_R(-\ba)$ で、クロスキャップ部分の差は閉じた形 $\pm\psi^2[\delta]$ に尽きる。
\item \textbf{正値性（§\ref{sec:positivity}）}：多箱では符号がそろわない。障害は、凍結した状態の多重項の基底と、$E$・$C$ の分離にある。$E$ と $C$ の分離は指数からは決まらず、composite の DAHA の値からも状態の数は得られなかった。framed unknot では $E$ が framing によらず一意に決まる。多箱の $c = 1$ の $\sigma$ 非対称性は、$c$ の偶奇の規則（hatted では $\beta \equiv |R| - 1$）と、refined に固有の $\sigma$ 奇部分に分かれる。
\end{itemize}

\begin{table}[t]
\centering
\small
\caption{本論文の検証の範囲。「整数」は整数係数の Laurent 多項式（$|R| = 2$ は $2g$ が整数で $g_{(2)} + g_{(1,1)}$ が整数）。「普遍因子」は §\ref{sec:universal} の割り切れ。}\label{tab:summary}
\begin{tabular}{p{3.6cm}p{2.6cm}p{2.4cm}p{5.4cm}}
\toprule
対象 & 整数 & 普遍因子 & 正値性（符号のそろい方） \\
\midrule
unknot、Hopf（$a, b \le 4$） & 成立 & 成立（$\hat C_{R,S}$ の 10 成分） & $\hat g_{R,[1]}$ はそろう。$\hat C_{R,S}$ の refined だけの部分は 8/10 \\
$T(2,2m+1)$ の [1]（$m \le 5$） & 成立 & $c = 2$ 部分は $m \le 5$、全体は $m \le 2$ & 判定式が成立（商の符号は $m \le 2$ で確認） \\
$T(2,2m+1)$ の固定 $N$（$N = 2, 3$、$|R| \le 3$、$m \le 3$） & 全方向の hatted で成立 & — & 固定 $N$ では判定できない \\
trefoil の 2 箱（大きな $N$） & 成立 & 成立 & 商のブロックの 5/10 \\
$T(2,5)$ の 2 箱（大きな $N$） & 成立 & 成立 & 6/18 \\
trefoil の 3 箱（大きな $N$） & 成立（$(1+t)$ の極なし、hatted も整数） & 成立 & 2/15, 3/14, 2/13 \\
refined Sp$(2N)$：$T(2,3)$, $T(2,5)$ の [1]、trefoil の 2 箱 & 成立（SO と同じ構造） & — & SO との関係 $g^{\mathrm{Sp}}(\ba) = -g^{\SO}(-\ba)$（[1]）、$W_R$ の $(-1)^{|R|}$（2 箱） \\
\bottomrule
\end{tabular}
\end{table}

\subsection{主張の分類}
本論文の主張は、解析的に導いたもの（定理・命題・補題）、計算機で確かめたもの（計算結果。範囲を明記する）、予想、仮説に分けて書く。
文献から直接取れなかった規則は「仮説」と書き、どの検証で裏付けたかを添える。
計算のコードとデータ、自動テストの対応は付録 \ref{app:code} に置く。

\section{準備}\label{sec:prelim}

\subsection{変数}
\begin{center}
\small
\begin{tabular}{>{$}l<{$}p{9cm}}
\toprule
\text{記号} & 意味 \\
\midrule
q,\ t & refined CS の変数。$r = q^{1/2}$, $s = t^{1/2}$ \\
U = t^N & SO$(2N)$ の Kähler 変数（$N$ は D$_N$ のランク）。大きな $N$ では形式変数 \\
\nu = U t^{-1/2} & Mariño の $\nu$（$t = q$ で $\nu = q^{N-1/2}$） \\
\tnu = \nu\, (q/t)^{1/2} & ずれた Kähler 変数 \\
A = \nu^2 = t^{N_A},\ N_A = 2N - 1 & 被覆 U$(N_A)$ の Kähler 変数 \\
(\bq, \bt, \ba) = (t^{1/2}, -(q/t)^{1/2}, \nu) & 太字の変数（DGR \cite{DGR05}・\cite{AgSc12} と同じ）。$q^{1/2} = -\bt\bq$ \\
V = (qt)^{1/4},\ W = (q/t)^{1/4} & $q^{1/2} = VW$, $t^{1/2} = V/W$ \\
z_q = q^{1/2} - q^{-1/2},\ z_t = t^{1/2} - t^{-1/2} & $z_q^+ = q^{1/2} + q^{-1/2}$ \\
G = z_q z_t & $G$ 展開の変数（KN \cite{KN17}） \\
\sigma\colon q \mapsto t^{-1},\ t \mapsto q^{-1} & $\nu$ を固定。$G$ は不変 \\
\bottomrule
\end{tabular}
\end{center}
unrefined の Mariño の変数は $t_M = q^{1/2}$、$z = z_q$ である。

\subsection{refined SO$(2N)$ の不変量}
AgSc \cite{AgSc12} の SO$(2N)$ の refined Chern--Simons 理論では、表現 $\lambda$ の Wilson ループは D$_N$ 型の Macdonald 多項式で与えられる。
% [S1]
\begin{compresult}
AgSc の refined unknot $\dim_\lambda$ は、Koornwinder 多項式 $K_\lambda(x; q, t; 1, -1, q^{1/2}, -q^{1/2})$ の主特殊化 $x = (1, t, \dots, t^{N-1})$ に等しい（6 つの $\lambda$。van Diejen--Sahi の評価式と、$q$-差分作用素で作った多項式（$N = 2, 3$）で確認）。
$\dim_\lambda$ は $U$ の Laurent 多項式ではなく、例えば [2] は $(1 - q t^{N-1})$、[1,1] は $(1 - q t^{2N-3})$ を分母にもつ（$t = q$ では約分される）。
\end{compresult}

% [S2]
\begin{remark}
Mariño 型の構成（bosonic な Cauchy 核、$\tfrac12 f^{\cov}$、$\sum_{d\ \mathrm{odd}}$）をそのまま refined にすると、unknot の 2 箱で $U$ の極が打ち消さない。
composite $([1],[1])$ の寄与は $p_1^2$ 方向だけにあるが、D$_N$ の [2] と [1,1] の極の位置は違う（$U = t/q$ と $U^2 = t^3/q$）。
重み 5 通りと基底 3 通りのどの組でもこれは解消しない。
\end{remark}

\subsection{双対 Cauchy 核と生成関数}
% [S3]
Mimachi の双対 Cauchy 恒等式（Koornwinder 多項式）
\[
\sum_{\lambda \subseteq m^n} (-1)^{|\lambda|} K_{m^n - \lambda}(x; q, t)\, K_{\lambda'}(y; t, q) = \prod_{i, j} x_i^{-1}(1 - x_i y_j)(1 - x_i y_j^{-1})
\]
では、ホロノミー側に AgSc の D$_n$ 型（$\pm1, \pm q^{1/2}$）、probe 側に $(q', t') = (t, q)$ に対する $(\pm 1, \pm t'^{1/2})$ 型が現れる。
この恒等式は $(n, m) = (1,1), (2,1), (1,2), (2,2), (2,3), (3,2)$ の 6 通りで確かめた。
これを refined orientifold の Ooguri--Vafa 核と読み、probe 変数 $Y = (y_1, y_2, \dots)$ に対して
\begin{equation}\label{eq:Z}
\tilde Z(Y) = \Bigl\langle \prod_{\alpha, i} (1 - y_\alpha x_i)(1 - y_\alpha x_i^{-1}) \Bigr\rangle_{\SO(2N)},\qquad
\log \tilde Z = \sum_{d \ge 1} \frac1d \sum_R f_R(q^d, t^d, U^d)\, s_R(Y^d)
\end{equation}
と置く。被覆 U$(N_A)$ についても、同じ核で $f^{\cov}_R$ を定める。
核の $|Y| \le 2$ の部分は、U$(N_A)$ の普通の表現と随伴 composite $(1;1)$ で書ける。3 次では composite $(2;1)$, $(1,1;1)$ も現れる。

\begin{definition}[クロスキャップ部分]
全 $d$ の定義で $g_R = f_R - \tfrac12 f^{\cov}_R$ と置く。
hatted 不変量は $\hat g = M(q^{-1})^{-1} g$ である。ここで $M_{RS} = \sum_\mu \chi_R(\mu)\chi_S(\mu) z_\mu^{-1}\prod_i (t_M^{\mu_i} - t_M^{-\mu_i})/(t_M - t_M^{-1})$（$t_M = q^{1/2}$）は Mariño の $M$ で、KN の $M(x)$ の $x = q^{-1}$ に等しい。
\end{definition}

% [S8]
\begin{remark}
reformulation に $M(t)$ や $M(t^{-1})$ を使うと、Hopf 絡み目で Laurent にならない。$M(q)$ は $R \leftrightarrow R^T$ を入れ替える。
KN の二つの枠（q 枠：Macdonald の Cauchy 核と $M(t)$、t 枠：双対 Cauchy 核と $M(q^{-1})$）の hatted 不変量は、2 箱で一致する（§\ref{sec:generalN}）。
\end{remark}

\subsection{一様な規格化}\label{sec:uniform}
入力の U$(n)$ の refined 不変量（KN の (N1)+(N2) で規格化したもの）は、行列式表現と荷電共役について $(q/t)$ の冪だけずれる。
そこで、[1] の superpolynomial の $\delta$ 次数 $\delta = \ba + \bq/2 - \bt$ を使い、
\begin{equation}\label{eq:uniform}
W_\mu = (t/q)^{\delta |\mu|/2}\, R_\mu\, P_\mu(t^\rho)
\end{equation}
と置く（$T(2,2m+1)$ では $\delta = m$）。

% [S5][S6]
\begin{compresult}\label{cr:uniform}
\begin{enumerate}
\item 規格化 \eqref{eq:uniform} では、行列式のずらし $\mu \to \mu + 1^n$ と荷電共役で $W$ が正確に不変になる。検査した範囲は trefoil の 16 組（U(2)〜U(5)）、$T(2,5)$ と $T(2,7)$ の各 13 組、二重ツイスト結び目 7 個である。
\item 固定 $N$（$N = 2, 3$、$|R| \le 3$）で framing の単項式を「$N$ に依存する極の打ち消し」と「$t = q$ で 1」で決めると、SO 側・被覆側ともすべて 1 になる（$T(2,3), T(2,5), T(2,7)$）。この走査では $(t/q)^{\beta/2}$ の族の中で $\beta$ が一意に決まる。
\item DAHA-Jones 多項式（鉛直 framing）とは、framing 固有値の冪を除いて一致する（§\ref{sec:generalN}）。
\item $t = q$ では Mariño の標準 framing そのものになる。
\end{enumerate}
\end{compresult}

% [S7]
\begin{hypothesis}
一様な規格化 \eqref{eq:uniform} は framing 0 の値である。
\end{hypothesis}
以後、SO 側・被覆側ともこの規格化を使う。

\section{orientifold の refined 指数と全 $d$ の構造}\label{sec:index}

\subsection{設定}
AgSc \cite{AgSc12} §3 にしたがい、M 理論を $X \times \mathrm{TN} \times S^1$ 上で $(q, t)$ で捻る。orientifold の作用 $\tilde I$ は、$X$ の反正則な対合 $I$ と、時空の $z_2 \to -z_2$ の組である。
結び目 $K$ の probe は M5 brane で、$z_2$ 平面を張る（t-brane。開弦の分母が $1/(1-t)$ であることによる）。被覆には $L_K$ と像 $I(L_K)$ の 2 枚がある。
被覆の一粒子状態の空間は $H_1 = H_{\mathrm{int}} \otimes F_{z_2}$ である。$F_{z_2}$ は $z_2$ 方向のモード $t^l$（$l \ge 0$）から成る。$\tilde I$ は $F_{z_2}$ に $(-1)^l$ で働く。

\begin{hypothesis}[一粒子状態の分類]\label{hyp:classes}
商の一粒子状態は被覆の $\tilde I$ の軌道で、AgSc の積表示 (3.13) のとおり次の 3 種に分かれる。
\begin{enumerate}
\item $\tilde I$ で入れ替わる対 $\{\psi, \tilde I\psi\}$：全モードの塔をもつ。
\item $\tilde I\psi = \varepsilon\psi$（$\varepsilon = \pm 1$）：$(-1)^l\varepsilon = 1$ のモードだけをもつ。
\item $z_2 = 0$ に凍結した状態：モードの塔をもたない。
\end{enumerate}
\end{hypothesis}

\subsection{全 $d$ の構造}
% [O1][O2]
\begin{proposition}\label{prop:alld}
仮説 \ref{hyp:classes} のもとで、状態の電荷を $m$ とすると
\[
g = f - \tfrac12 f^{\cov} = \frac12\,\frac{E}{1+t} + C,\qquad E = \sum_{\text{不変}} \varepsilon_\psi\, m_\psi,\qquad C = \sum_{\text{凍結}} m_\psi .
\]
\end{proposition}
\begin{proof}
被覆の指数は $f^{\cov} = \sum_{\text{対}} 2m/(1-t) + \sum_{\text{不変}} m/(1-t)$ である。
商の指数は $f = \sum_{\text{対}} m/(1-t) + \sum_{\text{不変}} m \sum_{(-1)^l\varepsilon = 1} t^l + \sum_{\text{凍結}} m$ である。
$\sum_{(-1)^l\varepsilon = 1} t^l = \tfrac12[1/(1-t) + \varepsilon/(1+t)]$ を代入すれば主張が従う。
\end{proof}

\begin{remark}\label{rem:Esector}
$\tilde I$ で不変な状態は、$\tilde I$ で自分自身に移るセクター（brane と image を同数もつセクター）にしか無い。したがって $E$ はそのセクターにだけ現れる（命題 \ref{prop:boundary} の系も参照）。
$\tilde I$ が brane と image の入れ替えとして働くなら、そのトレースは冪和の基底で $p_2$ 型の方向（Adams の $\psi^2$）に入り、$E_{(2)} = -E_{(1,1)}$ となる。
後半は Hopf 絡み目（計算結果 \ref{cr:E22}）と結び目の 2 箱（計算結果 \ref{cr:largeN}）で確かめた性質である。
\end{remark}

\subsection{trace 表示：Witten の二つの公式の SO$(2N)$ 版}
Witten \cite{Witten11} は、Chern--Simons の分配関数を 5 次元（6 次元 $(2,0)$ 理論）の $Q$ コホモロジー上のトレース $Z^{\mathrm{CS}} = \mathrm{Tr}_{\mathcal K}\, q^{P}(-1)^{F}$ と書き（§4.1）、単純レースでない群については、大きい理論の空間で外部自己同型 $\zeta$ を挿入したトレースと、固定された理論の普通のトレースという二つの表示を与えた（§3.4）。
D 型の $(2,0)$ 理論は、M5 brane を $X/\mathbb{Z}_2$（$\mathbb{Z}_2$ は法束に $-1$）に巻いて得られる（同論文 §5.1.4、\cite{AgSc12} 付録 A）。これが本論文の $\tilde I$ の 5-brane 側の起源である。
% [O9]
\begin{proposition}\label{prop:trace}
仮説 \ref{hyp:classes} のもとで、商の一粒子の指数は被覆の射影 $P_+ = (1 + \tilde I)/2$ と凍結した状態の和で書け、
\[
g = f - \tfrac12 f^{\cov} = \tfrac12\,\mathrm{Tr}_{H_1^{\cov}}\bigl(\tilde I\cdot(\text{指数の重み})\bigr) + C
\]
である。右辺の第 1 項は $\tfrac12 E/(1+t)$ に等しい。
\end{proposition}
\begin{proof}
$f = \mathrm{Tr}_{H_1^{\cov}}(P_+\cdot) + \sum_{\text{凍結}} m$ と、$\tilde I$ が不変な状態に $\varepsilon$、$z_2$ のモードに $(-1)^l$ で働くことから、$\mathrm{Tr}(\tilde I\cdot) = \sum_{\text{不変}}\varepsilon m\sum_l (-t)^l = E/(1+t)$。
\end{proof}
$E$ は Witten の「$\zeta$ を挿入したトレース」の型、$C$ は商の理論にしかない状態の寄与である。命題 \ref{prop:EC} の不定性は、一つの数を二つの表示に分ける仕方の不定性と読める。

\subsection{境界ゼロモードへの $\tilde I$ の作用}
KN \cite{KN17} §3.2 では、境界 $h$ 個の曲線の境界ゼロモード $d\theta_i$（$\sum_i d\theta_i = 0$）が hook 表現を張り、t-brane 枠では 1 モードあたり $(1 - q^{-1})$（規格化 $q^{1/2}$）を与える。
% [O10]
\begin{proposition}\label{prop:boundary}
KN の境界ゼロモードの描像を仮定し、$\tilde I$ が $L$ の境界と $I(L)$ の境界を入れ替えるとする。$\tilde I$ 不変な曲線は $L$ と $I(L)$ の境界を同数 $n$ 組もち、対ごとに和 $d\theta_i + d\theta_i'$ は偶、差 $d\theta_i - d\theta_i'$ は奇である。拘束 $\sum d\theta = 0$ は偶のモードを一つ減らす。したがって
\[
E\ \text{の境界因子} = z_q^{\,n-1}(z_q^+)^{\,n},\qquad
P := (1-t)f^{\cov}_{\text{固定セクター}}\ \text{の境界因子} = z_q^{\,2n-1}.
\]
系として、$E$ は brane と image の境界を同数もつセクター（$p_2$ 型の方向）にだけ現れる。
\end{proposition}
\begin{proof}
奇のモードでは、$\tilde I$ を挿入したトレースで $(1 - q^{-1})$ が $(1 + q^{-1})$ に替わる。偶 $n-1$ 個、奇 $n$ 個を掛け合わせればよい。
\end{proof}
% [O10]
\begin{compresult}\label{cr:boundary}
保存済みのデータで、命題 \ref{prop:boundary} の予言がすべて一致した。
\begin{center}
\begin{tabular}{llll}
\toprule
対象 & $n$ & $E$ の境界因子 & $P$ の境界因子 \\
\midrule
Hopf $(2,2)$, $(2,4)$ & 2 & $z_q(z_q^+)^2$ & $z_q^3$ \\
$T(2,3)$, $T(2,5)$ の $p_2$ 方向 & 1 & $z_q^+$（整数計画は $z_q^+$ でなければ解なし） & $z_q$ \\
framed unknot（framing 2, 4） & 1 & $z_q^+$ & $z_q$ \\
\bottomrule
\end{tabular}
\end{center}
\end{compresult}
計算結果 \ref{cr:E22} の $E_{22}$ の $(z_q^+)^2$ と、§\ref{sec:positivity} の「結び目の境界因子は和でなければならない」は、この命題の帰結である。残る未知は、モジュライ空間の部分 $H^*(\mathcal{M}_{g,h,\beta})$ への $\tilde I$ の作用（$\beta$ ごとの符号 $\varepsilon$）である。

\subsection{$\sigma$ 対称性の起源}
$\sigma\colon q\mapsto t^{-1}, t\mapsto q^{-1}$ は、KN の変数 $q_\ell = (qt)^{1/2}$, $q_r = (q/t)^{1/2}$ では $q_\ell \mapsto q_\ell^{-1}$、$q_r$ 固定、すなわち Jacobian の Lefschetz 作用 SU$(2)_\ell$ の Weyl 反射である。
KN の $\hat f$ の $\sigma$ 不変性は、取っ手のゼロモードが SU$(2)_\ell$ の二重項をなすこと（取っ手あたり $-(q_\ell + q_\ell^{-1}) + (q_r + q_r^{-1}) = G$）から来る。
凍結した状態は向き付けられない面 $N = \hat\Sigma/I$ に巻いた M2 で、$N$ には複素構造が無く $H^1(N)$ に Lefschetz の $\mathfrak{sl}_2$ は働かない。したがって $C$ に $\sigma$ 対称性を期待する理由は無く、多箱の $c = 1$ 部分が $G$ 展開できないこと（§\ref{sec:generalN}）はこれで説明できる（予想）。[1] の $c = 1$ 部分が $\sigma$ 対称であることは、次の観察で説明できる（予想）。[1] の $c = 1$ 部分の $\nu^\beta$ のブロックの $G$ の次数は、brane--image のセクター $(t^{-1/2}-t^{1/2})x/z_q$ の $\nu^{2\beta}$ のブロックの $G$ の次数のちょうど半分である（$T(2,3)$, $T(2,5)$ のすべての $\beta$ で確認）。[1] の凍結した状態は brane--image の環状の弦の商（Möbius の帯）で、商の取っ手 $g$ 個が被覆の取っ手 $2g$ 個に当たり、商の取っ手が $G$ を担う。多箱の $c = 1$ 部分が $\sigma$ 対称でない理由は未解明である。

Mariño の $\sum_{d\ \mathrm{odd}}$ の定義 $g^{\mathrm{odd}}$ とは、$g_{ab} = g^{\mathrm{odd}}_{ab} - \tfrac12 g^{\mathrm{odd}}_{a/2, b/2}(x^2)$ で結ばれる。

% [O4]
\begin{compresult}
Hopf 絡み目（$a, b \le 4$、§\ref{sec:hopf}）で、全 $d$ の $g_{ab}$ の分母は（偶, 偶）のセクターで $2(1+t)\times$ 単項式、それ以外で単項式である。
$\sum_{d\ \mathrm{odd}}$ の定義で現れる $4(1+t)(1+t^2)$ は、全 $d$ の定義では打ち消す。
\end{compresult}

\subsection{$E$ と $C$ の分離の不定性}
% [O3]
\begin{proposition}\label{prop:EC}
整数係数の Laurent 多項式 $X$ について、$E \to E + 2(1+t)X$, $C \to C - X$ としても $g$ は変わらない。したがって $g$ から決まるのは $E$ の $2(1+t)$ を法とした類だけで、特に $E$ そのものは決まらない（$E(-1) = [2(1+t)g]_{t=-1}$ は決まる）。
\end{proposition}
\begin{proof}
$\tfrac12 \cdot 2(1+t)X/(1+t) = X$ なので、$E$ の変化は $C$ の変化で打ち消す。
状態の言葉では、次の二つの状態の組は、商の指数 $f$ と被覆の指数 $f^{\cov}$ の両方で区別できない。
\begin{itemize}
\item 凍結した状態 $m$ と、$\tilde I$ で入れ替わる対 $\{m\}$, $\{mt\}$。
\item $\varepsilon = +1$ の不変な状態 $m$, $m$, $mt$, $mt$。
\end{itemize}
実際、$f^{\cov}$ への寄与はどちらも $2(m + mt)/(1-t)$、$f$ への寄与はどちらも $2m/(1-t)$ である。
\end{proof}

% [O7]
\begin{corollary}\label{cor:Efree}
冪和の $p_1^n$ 方向と、箱の数が奇数のセクター（$E = 0$）は、$E$ によらない。これらの方向では、正値性を $E$ の不定性によらずに判定できる。
\end{corollary}

\subsection{閉弦での確認}
% [O5]
\begin{compresult}\label{cr:closed}
D$_N$ の積表示 $Z = \prod_{m<\beta}\prod_{\alpha>0}(q^{-m/2}t^{-(\alpha,\rho)/2} - q^{m/2}t^{(\alpha,\rho)/2})$（$t = q^\beta$）から作った $-\log Z$ の $d$ 番目と、AgSc §6 の 3 項
\[
\frac{1}{2d}\frac{Q^d t^{d/2} q^{-d/2}}{z_q(q^d) z_t(t^d)} + \frac1d \frac{(Q t/q)^{d/2}}{z_q(q^d)} - \frac{1}{2d}\frac{Q^d t^{d/2}q^{-d/2}}{z_q(q^d)(t^{d/2} + t^{-d/2})},\qquad Q = t^{2N-1}
\]
との差は、$N$ の 1 次式（古典的な項と定数写像）だけである（$N = 2, \dots, 7$）。
\end{compresult}
AgSc の 3 項は、命題 \ref{prop:alld} の閉弦版である。P$^1$ は $\varepsilon = +1$ の不変な状態で、第 1 項と第 3 項の和が $z_2$ の偶数のモードだけの寄与になる。RP$^2$ は凍結した状態で、第 2 項に当たる。

Avetisyan--Mkrtchyan \cite{AM22} は、同じ分配関数を多重サイン関数で書き、D・C 型の自由エネルギーを 3 種に整理した。奇数の $n$ の $c = 1$ の項、全 $n$ の「Klein 瓶」（$c = 2$）の項、Kähler 変数の「三重化」である。
% [O6]
\begin{proposition}\label{prop:AM}
AM の変数 $(Q, T) = (q^{1/2}, t^{1/2})$ で、$X = e^{\tau_1/2}$, $e^{\tau_2} = X^2/Q$, $e^{\tau_0} = X^2 T/Q$ とする。このとき $n$ ごとの有理式として、次の恒等式が成り立つ。
\[
\underbrace{-\frac{X^{2n}}{2n(Q^{2n} - Q^{-2n})}}_{\text{奇数 $n$ に直す差}}
+ \underbrace{\frac{e^{n\tau_2}}{2n}\frac{Q^n T^{-n} - T^n Q^{-n}}{(Q^{2n} - Q^{-2n})(T^n + T^{-n})}}_{\text{Klein 瓶}}
= -\frac{e^{n\tau_0}}{2n\,(Q^n - Q^{-n})(T^n + T^{-n})} .
\]
したがって AM の「奇数 $n$ の $c = 1$ ＋ Klein 瓶」は、全 $n$ の凍結した状態の項 $e^{n\tau_1/2}/(n(Q^n - Q^{-n}))$ と、$\varepsilon = +1$ の不変な状態の項の和に等しい。
Klein 瓶の項とその分子 $Q^n/T^n - T^n/Q^n$（AM の論文の変数では $q/t - t/q$）は、独立な多重項ではない。奇数の $d$ で整理したときの $-\tfrac12\psi^2 C$ と $\tfrac12 E/(1+t)$ の差である（$E \ne \psi^2 C$）。
\end{proposition}
\begin{proof}
左辺を通分すると、$X^{2n}$ の係数は $[-(T^n + T^{-n}) + Q^{-n}(Q^nT^{-n} - T^nQ^{-n})]/(\cdots) = -T^nQ^{-n}(Q^n + Q^{-n})/(\cdots)$ になり、$(Q^{2n} - Q^{-2n}) = (Q^n - Q^{-n})(Q^n + Q^{-n})$ で約分できる。
\end{proof}

\begin{remark}
AM と AgSc の規約は、$(q, t, Q) \to (q^{-1}, t^{-1}, Q^{-1})$ と全体の符号（$\log Z$ と $-\log Z$）で移り合う。凍結した状態と向き付けられた状態の電荷の比は、AgSc で $x^2/Q_0 = w^{-1}$、AM で $w$ である（$w = (q/t)^{1/2}$）。
閉弦では状態が P$^1$ と RP$^2$ の二つだけで打ち消しが無いので、$E$ と $C$ は一意に決まる。命題 \ref{prop:EC} の不定性は、開弦で初めて問題になる。
\end{remark}

% [O8]
\begin{remark}[取っ手の因子]
AgSc §4 の強い整数性の導出に現れる「差 $\times$ 和」$(q^{1/2} - q^{-1/2})(q^{1/2} + q^{-1/2})$ は、類 $\beta/2$ の $c = 1$ の寄与を $d$ を倍にして並べたときの $\psi^2 z$ である。これは取っ手一つの因子ではない。
Piazzalunga--Uranga \cite{PU14} の M 理論の解釈では、$\tilde I$ で不変な被覆の取っ手一つの捻ったトレースは、同じ変数で $z$ 一つになる。
refined では $G = z_q z_t = T_+T_-$、$T_\pm = (V - V^{-1}) \pm (W - W^{-1})$ と分解できる。ただし $T_\pm$ を使っても、$\sigma$ 対称でない $c = 1$ 部分は展開できるようにならない。
$\mathbb{Z}[W^{\pm 2}][T_+, T_-] = \mathbb{Z}[W^{\pm2}][V - V^{-1}, W - W^{-1}]$ の元のうち $V$ の偶数冪の部分は、$(V - V^{-1})$ の偶数冪の和なので、$V \to V^{-1}$（$\sigma$）で不変だからである。
\end{remark}

\section{unknot と Hopf 絡み目}\label{sec:hopf}

\subsection{virtual 正規化}
probe 側の基底 $K_{(n)}(y; t, q; 1, -1, q^{1/2}, -q^{1/2})$（monic な Askey--Wilson 多項式）の $y^{n-2j}$ の係数は、$U_n = t^n$ の有理式として
\[
c_{2j}(n) = t^j\, \frac{(q/t; t^2)_j}{(t^2; t^2)_j}\cdot\frac{(1 - U_n)(1 + U_n t^{-2j})\prod_{i=1}^{j-1}(1 - U_n^2 t^{-2i})}{\prod_{i=1}^{j}(1 - qU_n^2 t^{-2i-1})}
\]
と書ける。
% [S4]
これは予想で、有理点 4 つで $j \le 5$、$n \le 30$ について漸化式と一致することを確かめた。
virtual 正規化 $\tilde Z = y^N z^N Z$ は $y, z$ の多項式で、$y^a z^b$ の係数は $N$ によらない有限和になる。したがって $U$ を形式変数のまま扱える（$U = t^N$、$N = 3, \dots, 8$ の直接計算と $a + b \le 8$ で全係数一致）。

\subsection{unknot}
% [H1]
\begin{proposition}
unknot では $f_{10} = (1 - U)(U + t)/(U(t-1))$、$f_{a0} = 0$（$a \ge 2$）である。$\nu$ で書くと
\[
f_{10} = -1 - \frac{\nu - \nu^{-1}}{t^{1/2} - t^{-1/2}}
\]
で、Mariño 型の unknot $(\nu - \nu^{-1})/z$ とクロスキャップの $-1$（Mariño の $g_{[1]} = 1$）の和である。積表示は $\prod_j (y_j^{\pm}; t)_\infty / (U y_j^{\pm}; t)_\infty$ で、$\tfrac12$ も $\sum_{d\ \mathrm{odd}}$ も要らない。
\end{proposition}
\begin{proof}
双対 Cauchy 核の期待値は $\prod_j\prod_{k=1}^{N} t^{-(N-k)}(1 - t^{N-k}y_j^{\pm})$ で、$\log$ をとればよい。
\end{proof}

\subsection{Hopf 絡み目}
二つの成分に probe を一つずつ付け、$f_{ab} = \sum_{d | (a,b)} (\mu(d)/d)\, L_{a/d, b/d}(q^d, t^d, U^d)$ と置く。
% [H2][H3]
\begin{compresult}\label{cr:hopf}
$a, b \le 4$ で次が成り立つ。
\begin{enumerate}
\item $f_{ab} \in \dfrac{1 - q}{1 - t}\,\mathbb{Z}[q^\pm, t^\pm, U^\pm]$。$(a, b)$ がともに偶数のときは分母が $1 - t^2$ になる。
\item 成分 1 の probe を無限個にした Schur 基底の $f_{R,b}$（$|R| \le 4$、$b \le 2$）でも同じ整数性が成り立つ。素朴な Macdonald 基底では $(1 - qt)$ などの余分な分母が出る。
\item $f_{11} = -f_{10}\, z_q(\tnu - \tnu^{-1})$、$f_{1b} = f_{11}\,\chi_{b-1}(U/t)$（$\chi_m(x) = x^m + x^{m-2} + \dots + x^{-m}$）。$t = q$ では Mariño の $\hat g_{[1],[1]} = z(\nu - \nu^{-1})$ に、符号と unknot の因子を除いて一致する。
\end{enumerate}
\end{compresult}
refined では、Kähler 変数が $\nu \to \tnu = \nu (q/t)^{1/2}$ とずれる。これは AM の Kähler 変数の量子シフト（§\ref{sec:index}）の開弦版と見なせる。

\subsection{クロスキャップ部分}
被覆には、U$(N)$ の refined CS の composite 表現での Hopf 絡み目を使う。$\tfrac12 f^{\cov}_{10} = -(\nu - \nu^{-1})/z_t$、$\tfrac12 f^{\cov}_{11} = (\nu - \nu^{-1})(\tnu - \tnu^{-1})z_q/z_t$ である。
% [H4][H5]
\begin{compresult}\label{cr:hopfg}
\begin{enumerate}
\item $g_{1b} = z_q(\tnu - \tnu^{-1})\,\chi_{b-1}(\tnu/q^{1/2})$（$b \le 4$）。$t$ は $\tnu$ を通してしか入らない。
\item $\hat g_{R,[1]}$（$|R| \le 4$）は整数 Laurent である。hook $R = (n-d, 1^d)$ では $(-1)^d z_q(\tnu - \tnu^{-1})\tnu^{2d-n+1}$、hook でない $R$ では 0 になる。
$t = q$ では Mariño の $\hat g_{[2],[1]} = z(\nu^2 - 1)$, $\hat g_{[1,1],[1]} = z(\nu^{-2} - 1)$ に、$R \leftrightarrow R^T$ と全体の符号 $-1$ で一致する。
\end{enumerate}
\end{compresult}

\subsection{$E$ の決定と Adams}
probe を $x$ 側と $1/x$ 側に分けたセクター別の被覆から、固定セクター（$\tilde I$ で自分に移るセクター）の状態の和 $P$ を作る。
$E = P - 2Q$（$Q$ は $\tilde I$ の固有値 $-1$ の状態の和、各単項式で $0 \le Q_m/P_m \le 1$）と書き、次の 2 条件を整数計画で課す。
(i) $t = -1$ で $E = 2(1+t)g$。
(ii) $t = q$ で $C$ が Mariño の値。
% [H6][H7]
\begin{compresult}\label{cr:E22}
\begin{align*}
E_{22} &= -z_q (z_q^+)^2(\tnu - \tnu^{-1})(t q^{-1/2}\tnu + q^{1/2}\tnu^{-1})\quad(\text{一意}),\\
E_{24} &= E_{22}\,(\tnu^2/q + q/\tnu^2)\quad(\text{一意}).
\end{align*}
$(4,4)$ は解が存在するが一意でない（固定セクターの単項式 180 個のうち 164 個が自由）。
成分 1 を多 probe にすると、$\tilde I$ は「固定セクターの入れ替え＋Adams」として働き、$E_{R,S} = \sum_{\rho, \sigma\vdash 2} c^R_{2;\rho}c^S_{2;\sigma}\,\varepsilon_{\rho\sigma}$ が成り立つ。確かめた範囲は、$|R| \le 4$, $|S| \le 2$ の 10 セクターと、$|R|, |S| \le 4$ の $t = -1$ での値 25 組である。
\end{compresult}
$P_{22} = z_q^3(\tnu - \tnu^{-1})(tq^{-1/2}\tnu - q^{1/2}\tnu^{-1})$ と比べると、$E_{22}$ は差の因子 $z_q^2$ を和の因子 $-(z_q^+)^2$ に替え、第 2 因子の符号を反転したものである。
これは $\tilde I$ の捻ったトレースの形で、§\ref{sec:positivity} で結び目に当てる。

% [H8]
\begin{compresult}\label{cr:hopfC}
成分 1 を多 probe（$|R| \le 4$）、成分 2 を 2 箱（$|S| \le 2$）にした $\hat C_{R,S} = M^{-1}(q^{-1})\,(g - \tfrac12E/(1+t))\,M^{-1}(q^{-1})^T$ を、$t = q$ の部分（$\tnu$ で表す）と refined だけの部分に分ける。
\begin{enumerate}
\item refined だけの部分は、10 成分すべてで $(\tnu^2 - q)(q - t)$ で割り切れる。太字では $\bt^2(\ba^2 - \bq^2)\cdot\bq^2(\bt^2 - 1)$ である。
\item $t = q$ の部分は 10/10 で、$\tnu$ の冪ごとに符号がそろう。
\item refined だけの部分の商は 8/10 でそろう。転置の対 $\hat C_{(3,1),(1,1)}$, $\hat C_{(2,1,1),(2)}$ では、$\tnu^2$ のブロックに $q^{1/2}(1-t)$ が入ってそろわない。
\end{enumerate}
\end{compresult}

\section{トーラス結び目の [1]}\label{sec:torus}

\subsection{入力と h/g 分解}
入力は AgSc \cite{AgSc12} §6 の $T(2,2m+1)$ の [1] の一般式（unknot で正規化）である。
$m = 1$ で AgSc (6.15) に一致する。
SO(4) $=$ SU(2)$^2$、SO(6) $=$ SU(4) の同一視で、(6.15) は Gorsky--Negut \cite{GN13} の構成による U$(n)$ の値と一致する（有理点 2 つで厳密）。
$t = q$ では Mariño の Kauffman 多項式の鏡像に、framing の単項式を除いて一致する。

% [T1]
\begin{hypothesis}[h/g 分解]\label{hyp:hg}
$f_1(K) = -W_{\SO}(K) = \tfrac12 h_{[1]} + g_{[1]}$ と書き、被覆の部分を
\[
\tfrac12 h_{[1]} = \mathcal{P}(\ba_{\mathrm{KN}}, \bq, \bt)\,\frac{\lambda - \lambda^{-1}}{\bq - \bq^{-1}},\qquad \ba_{\mathrm{KN}}^2 = a\,(t/q)^{1/2}
\]
と置く。$\mathcal P$ は DGR \cite{DGR05} の superpolynomial で、A 型の refined CS の [1] に当たる。
\end{hypothesis}

% [T2]
\begin{compresult}
$m = 1, \dots, 5$ で、$g_{[1]}$ が Laurent になる単項式 $\pm q^{c/2} t^{d/2} (t^N)^e$ を走査した。条件を満たすのは $e = -2(2m+1)$、$c = -(3m+2)$ に限られる。
$d$ は「$t = q$ の極限が Kauffman 多項式の形」で一意に $d = 7m+4$ と決まる。
したがって単項式は $T_{[1]}^{-2(2m+1)}(q/t)^{\delta/2}$ である（$T_{[1]} = q^{1/2}t^{N-1}$、$\delta = m$）。
一様な規格化では、余分な $(q/t)^{\delta/2}$ は被覆側に移り、SO 側の単項式はちょうど $T_{[1]}^{-2(2m+1)}$ になる（固定 $N = 2, 3$ で data と全係数一致、$m \le 3$）。
\end{compresult}
同じ形の因子 $\sqrt{q/t}$ は、Aganagic--Shakirov \cite{AgSh11} の SU$(N)$ の trefoil でも現れる。

\subsection{結果}
一様な規格化でのクロスキャップ部分を $g^{\mathrm{true}}_{[1]} = (t/q)^{\delta/2} g_{[1]}$ と書く。
% [T3][T4]
\begin{compresult}\label{cr:torus}
$m = 1, \dots, 5$ で、$g^{\mathrm{true}}_{[1]}$ は整数係数の Laurent 多項式（項数 11, 27, 49, 77, 111）で、次が成り立つ。
\begin{enumerate}
\item[(N)] SO(2) の点 $\ba^2 = \bq^2$（$\lambda^2 = t$）で $g^{\mathrm{true}}_{[1]} = 1$。
\item[(P)] $\bt \to -\bt$ で不変。
\item[($c{=}2$)] $\ba$ の奇数冪の部分は
\[
-(\ba^2 - \bq^2)(\bt^2 - 1)\,\ba^{-(4m+1)}\bq^{-(2m-1)}\bt^{-4m}\,\Pi_m,\qquad
\Pi_m = \sum_{j=0}^{m-1}\ba^{2j}\bq^{2(m-1-j)}\bt^{2j}\sum_{i=0}^{j}(\bq^4\bt^2)^i .
\]
$\Pi_m$ は係数がすべて $+1$ の $m(m+1)/2$ 項から成る。閉じた形は $m \le 3$ から推測し、$m = 4, 5$ で確かめた（予想）。
\item[($c{=}1$)] $\ql$ の整数冪の部分（$\ql^{1/2}$ の偶数冪）は $\sigma$ で対称で、$G$ と $\qr^{1/2}$ の多項式に書ける。$\lambda^\beta$ ごとのブロックは次のどれかである。
  \begin{itemize}
  \item $\lambda^{-(4m+2)}$：正の単項式
  \item $\lambda^{-4m}, \dots, \lambda^{-(2m+4)}$：$(\bt^2 - 1)\times$ 正
  \item $\lambda^{-(2m+2)}$：係数がすべて負
  \item $\lambda^{-2m}$：係数がすべて正
  \end{itemize}
\end{enumerate}
\end{compresult}
$t = q$ では $z$ の奇数冪が消え、トーラス結び目の [1] で $N^{c=2} = 0$ になることと整合する。
$m = 1, 2, 3$ の値は付録 \ref{app:torus} に載せた。

\subsection{判定式}\label{sec:criterion}
unrefined のクロスキャップの BPS 数は整数だが、符号は $\beta$ によって変わる（例：Mariño の Hopf の $\hat g_{[1],[1]} = z(\lambda - \lambda^{-1})$）。
refined では正値性を期待し、それを「各 $\lambda^\beta$ のブロックの係数が、unrefined の符号 $\varepsilon_\beta$ と同じ符号にそろう」こととして定式化する。

\begin{definition}\label{def:c1}
{[1]} のクロスキャップ部分の $c = 1$ 部分を
\[
g^{c=1} = \sum_\beta \lambda^\beta \sum_{g, J} N^{c=1}_{g,\beta,J}\, G^g \qr^{2J}
\]
と展開する。$t = q$ で $\sum_J N^{c=1}_{g,\beta,J}$ は Mariño の $N^{c=1}_{g,\beta}$ に戻る。
\end{definition}

% [T6]
\begin{conjecture}[判定式]\label{conj:criterion}
結び目の [1] では、次が成り立つ。
\begin{enumerate}
\item[($c{=}1$)] $G$ 展開の後、$\lambda^\beta$ ごとのブロックは同符号か、$(\bt^2 - 1)\times$ 同符号になる。
\item[($c{=}2$)] $c = 2$ 部分は $-(\ba^2 - \bq^2)(\bt^2 - 1)\times$ 正の多項式である。
\end{enumerate}
(N)(P) は規格化の検査で、正値性の内容は ($c{=}1$)($c{=}2$) にある。
\end{conjecture}
根拠はトーラス結び目 $T(2,2m+1)$（$m \le 5$）の [1] だけである。
$(\ba^2 - \bq^2)$ は SO(2) の点（可換群では結び目の不変量が自明になる）で消える因子である。
$(\bt^2 - 1)$ は「$t = q$ で消える」ことと、(P) の偶奇の帰結である。

\section{一般の $N$ の多箱}\label{sec:generalN}

\subsection{DAHA-Jones 多項式}
refined SO$(2N)$ の多箱の値は、AgSc の SL$(2,\mathbb{Z})$ の計算では一般の $N$ で得られない。そこで Cherednik \cite{Cherednik11} の DAHA-Jones 多項式
\[
JD_{r,s}(b) = \{\hat\gamma_{r,s}(P_b)\}/P_b(q^{-\rho_k})
\]
を、一般の $t$ で D$_N$ と GL$_N$ の多項式表現の上で直接計算する（同論文の shift 作用素の方法は $t = 1, q$ に限る）。

規約は次のとおりである。
\begin{itemize}
\item $(T_i - t^{1/2})(T_i + t^{-1/2}) = 0$。評価は $\{H\} = (H\cdot 1)(q^{-\rho_k})$。
\item ベクトル表現は minuscule なので $Y_{\varepsilon_1} = \pi_1 T_{u_1}$ で、$\pi_1(X_b) = q^{(\omega_{1*}, b)}X_{u_1^{-1}b}$（D$_N$：$\omega_{1*} = \varepsilon_1$、GL$_N$：$\omega_{1*} \equiv -\varepsilon_N$）。
\item $\tau_-(X) = q^{1/2}YX$、$\tau_+(Y) = q^{-1/2}XY$。$T(2m+1, 2)$ は $\gamma = \tau_+^m\tau_-^2$。
\end{itemize}
計算の高速化には、反自己同型 $\varphi$ で $X$ と $Y$ を入れ替え、軌道の対称化と $R$ の間の共有を使った。
被覆の composite $(\alpha;1)$ は、$X^{-1}$ 側の対称化で作った $\hat\gamma'(p_1(Y))\cdot 1$ を $X$ 側の連鎖の出発点にして計算した。

% [D1][D2][D3]
\begin{compresult}[検証]\label{cr:daha}
\begin{enumerate}
\item D$_N$, GL$_N$（$N = 2, 3, 4$）で、DAHA の関係式（2 次、組紐、交差、$Y$ の可換性）が成り立つ。D$_N$ の Macdonald 多項式は Macdonald 作用素の固有関数である。
\item D$_N$：$JD = T_R^{-2n}\times$ 鏡像（一様な規格化の固定 $N$ の $W_R$）。$T_R = q^{(R,R)/2}t^{(R,\rho)}$、$n = 2m+1$ で、確かめた範囲は $N = 2, 3$、$R \in \{(1), (1,1), (2)\}$、$m = 1, 2$ である。$N = 4, 5$ の [1] は AgSc (6.15) の鏡像にちょうど一致する。
\item GL$_N$：$W_b = $ 鏡像$(JD_b/M_b)$、$M_b = q^{-n(b,b)}t^{-2n(b,\rho) + (m+1)|b|(N-1)}$。確かめた範囲は次のとおりである。$m = 1$ では [1], [2], [1,1] から読み取り、[3], [2,1], [1,1,1] と随伴（U(3), U(5)）で予言どおりになった。$m = 2$ では $N = 2, 3, 4$ の [1], [2], [1,1]（9 例）と随伴（U(3), U(5)）で成り立つ。
\item 高速版と現行版は、保存した値で全係数一致する。
\end{enumerate}
\end{compresult}

\subsection{一般の $N$ の補間}
簡約した不変量 $\Rred_R(N) = T_R^{-2n}\times$ 鏡像$(JD(N))/K_R(t^\rho)$ を、$t^N$ の Laurent 多項式として補間した。
補間の窓は推測しない。$\Rred_R(N)$ の $t^{1/2}$ の次数の最小値・最大値が大きな $N$ で $2Nj + $（定数）になることから窓を決め、大きな $N$ をノードにして Newton の差分商で補間し、残りの $N$ で照合した。
% [D4]
\begin{compresult}\label{cr:interp}
次の表の範囲で、補間に使わない $N$ の値を全係数で予言した。
\begin{center}
\footnotesize
\begin{tabular}{llll}
\toprule
結び目・表現 & 冪の範囲（SO は $t^N$、被覆は $\lambda = t^{N_A/2}$） & ノード & 照合 \\
\midrule
trefoil の [1] & $-5$〜$-2$ & $N = 2$〜5 & $N = 6, 7$、AgSc (6.15) \\
trefoil の (2), (1,1) & $-10$〜$-4$ & $N = 2$〜8 & $N = 9$（$N = 3$〜9 から $N = 2$） \\
trefoil の (3), (2,1), (1,1,1) & $-15$〜$-6$ & $N = 4$〜13 & $N = 2, 3$（(1,1,1) は $N = 3$） \\
$T(2,5)$ の [1] & $-9$〜$-4$ & $N = 9$〜14 & $N = 2$〜8 \\
$T(2,5)$ の (2), (1,1) & $-18$〜$-8$ & $N = 4$〜14 & $N = 2, 3$ \\
被覆の随伴 $(1;1)$（trefoil） & $-10$〜$-4$ & $N_A = 2$〜11 のうち 7 点 & 残りの 3 点 \\
被覆の随伴 $(1;1)$（$T(2,5)$） & $-18$〜$-8$ & $N_A = 14$〜24 & $N_A = 2$〜13 \\
被覆の $(2;1), (1,1;1)$（trefoil） & $-14$〜$-6$ & $N_A = 12$〜20 & $N_A = 2$（$(1,1;1)$ は 3）〜11 \\
\bottomrule
\end{tabular}
\end{center}
3 箱を $t = q$ にすると、[3] は NRZ \cite{NRZ13} の [3] 彩色の Kauffman 多項式の鏡像に比 1 で一致する。
\end{compresult}
これで、refined SO$(2N)$ の trefoil の 2 箱・3 箱と $T(2,5)$ の 2 箱の不変量が一般の $N$ で得られた。
trefoil の 2 箱の係数は付録 \ref{app:daha} に載せた。

\subsection{大きな $N$ のクロスキャップ部分}
3 箱の SO 側は、核を D$_N$ の $K_R$（$|R| \le 3$、$t^N$ は形式変数）の逆行列で $K$ 基底に直して求めた。被覆は、普通の表現・随伴・composite から組み立てた。
% [D5][D6]
\begin{compresult}\label{cr:largeN}
\begin{enumerate}
\item 2 箱（trefoil, $T(2,5)$）：$g_{(2)}, g_{(1,1)}$ の分母は $2\times$ 単項式 $\times(1+t)$ だけで、$N$ に依存する極はすべて消える。$g_{(2)} + g_{(1,1)}$ と $(1+t)(g_{(2)} - g_{(1,1)})$ は整数 Laurent である。$(1+t)$ は $p_2$ 方向だけにあり、命題 \ref{prop:alld} の $E_{(2)} = -E_{(1,1)}$ と一致する。
\item 3 箱（trefoil）：$g_{(3)}, g_{(2,1)}, g_{(1,1,1)}$ は整数 Laurent で、$(1+t)$ の極は無い。奇数のセクターに $E$ が無いという命題 \ref{prop:alld} の予言どおりである。hatted の $\hat g_R = M(q^{-1})^{-1} g_R$ も $U$ 一般で整数 Laurent である（3 箱では全 $d$ の $g$ が Mariño の奇数 $d$ の $g^{\mathrm{odd}}$ に等しい）。
\item 検算として、$t^N = t^2, t^3$ で固定 $N$ の値と全係数一致する。
\end{enumerate}
\end{compresult}

% [D8]
KN の q 枠と t 枠の hatted 不変量は、2 箱の SO・被覆・クロスキャップ部分のすべてで一致する。

\subsection{$t = q$ での Mariño との辞書}
% [D7]
\begin{compresult}\label{cr:marino}
\begin{enumerate}
\item SO の $W_R$ と被覆の $W_S$ は、Mariño の $G_R, H_S$（標準 framing）の鏡像（$\lambda \to 1/\lambda$, $t_M \to 1/t_M$）に比 1 で一致する（$R, \mu \in \{(1), (2), (1,1)\}$）。
\item クロスキャップ部分では $g_{(1)} = -$鏡像$(g^{\mathrm{Mar}}_{[1]})$ である。2 箱は次の手順で $-$鏡像$(\hat g^{\mathrm{Mar}}_R)$ になる。
  \begin{itemize}
  \item 全 $d$ の $g$ を奇数 $d$ に直す：$g^{\mathrm{odd}}_{(2)} = g_{(2)} + \tfrac12 g_{(1)}(x^2)$、$g^{\mathrm{odd}}_{(1,1)} = g_{(1,1)} - \tfrac12 g_{(1)}(x^2)$。
  \item unknot の値 $\mp\tfrac12$ を引く。
  \item $M(q^{-1})^{-1}$ で hatted にする。
  \end{itemize}
\item Mariño \cite{Marino09} の (4.8) の $\hat g_{[2]}$ には誤植がある（$z^9$ の項）。正しくは $z^9(11\lambda^{12} - 78\lambda^{10} + 78\lambda^8 - 11\lambda^6)$ である。
\end{enumerate}
\end{compresult}
全体の符号 $-1$ は、双対 Cauchy（fermionic）の核の帰結である。$R \leftrightarrow R^T$ は奇数 $d$ に直した $g^{\mathrm{odd}}$ の段階で現れ（$g^{\mathrm{odd}}_{(2)}$ が鏡像$(g^{\mathrm{Mar}}_{[1,1]})$ に対応する）、鏡像で $M(q^{-1}) \leftrightarrow M(q)$ となることで hatted の段階では打ち消す。

\subsection{$G$ 展開の破れ}
% [D10]
\begin{compresult}
{[1]} の $c = 1$ 部分は $\sigma$ で対称だが、2 箱・3 箱では $c = 1$ 部分が対称にならない。対象は $\hat g_R$、和、奇数 $d$ 版のどれでも同じで、$c = 1$ を $\lambda$ の偶数冪ととっても、hatted の正しい偶奇（$\beta \equiv |R| - 1$、計算結果 \ref{cr:multibox-sigma}）でとっても同じである。
SO の簡約した不変量 $\Rred_R$ には、$\sigma$・$q\leftrightarrow t$・鏡像と、$\lambda \to \pm\lambda^{\pm1}\ql^{c_1}\qr^{c_2}$（$|c_1|, |c_2| \le 2$）、$R \leftrightarrow R^T$ をどう組み合わせても、比が単項式になる対称性が無い。
\end{compresult}
A 型の super-duality に当たる対称性が SO 側には無いので、SO の $\hat f$ と $\hat g$ が $\sigma$ 対称になる理由は無い（被覆の $\hat f$ は $R$ ごとに $\sigma$ 対称である）。

\subsection{refined Sp$(2N)$ の [1] と 2 箱}\label{sec:sp}
Avetisyan--Mkrtchyan \cite{AM22} の refined Chern--Simons の定義は、任意のルート系で $k_\alpha = y$ を一様にとり $q_\alpha = q^{(\alpha,\alpha)/2}$ とするもので、C$_N$ では長い根のパラメータが $u = t_{\mathrm{lng}} = t^2$ になる。
Cherednik \cite{Cherednik11} §Kauffman の C$_n$ の DAHA（$t = t_{\mathrm{sht}}$, $u = t_{\mathrm{lng}}$）を多項式表現の上で実装し、次を確かめた。
% [D11]
\begin{compresult}
\begin{enumerate}
\item C$_2$, C$_3$ で DAHA の関係式が成り立つ。$u = 1$ で D$_N$ の DAHA-Jones の [1] に比 1 で一致する（$N = 2, 3$、$T(3,2)$, $T(5,2)$）。Cherednik の $\mathcal{H}^C_{3,2}(\omega_1; q, t, u, a_{\mathrm{C}})$ に $a_{\mathrm{C}} = -t^{m-1}$ で単項式を除いて一致する（$m = 2, 3, 4$）。
\item $N = 2, \dots, 8$ の値から hyper-polynomial $\mathcal{H}^C_{2m+1,2}(\omega_1)$ を $a_{\mathrm{C}}$（Cherednik の $a$）の多項式として復元した。$T(5,2)$ の式は新しい。
\item SO の規約：$W_{\SO} = $ 鏡像$[\mathcal{H}^C(u{=}1, a_{\mathrm{C}}{=}-\lambda/t^{1/2})]\cdot\dim_V^{\SO}(\lambda)\cdot(t^{1/2}/\lambda)^{2m}$ で、$g^{\mathrm{true}}_{[1]} = W_{\SO} - W^{\cov}_1$ が再現される（一意）。
\item Sp：$\lambda = t^{N+1/2}$（$a = \lambda^2 = t^{2N+1}$。被覆は SO と同じ $\lambda$ の関数で取る）、$a_{\mathrm{C}} = -t^{N-1}$、$u = t^2$。単項式は $t = 1$ の極の打ち消しと $t = q$ での Kauffman の形から SO と同じ $(t^{1/2}/\lambda)^{2m}$ に決まり、$T(2,3)$, $T(2,5)$ で
\[
g^{\mathrm{Sp}}_{[1]}(\ba, \bq, \bt) = -\,g^{\SO}_{[1]}(-\ba, \bq, \bt)
\]
が全係数で成り立つ。
\end{enumerate}
\end{compresult}
$c = 1$ を $\ba$ の偶数冪、$c = 2$ を奇数冪に割り当てているので、これは $C^{\mathrm{Sp}} = -C_{c=1} + C_{c=2}$、すなわち Mariño の「SO と Sp は $c = 1$ の項の符号が違う」が refined でもそのまま成り立つことを意味する。$c$ の割り当てを独立に裏付けるとともに、AM の refined Sp$(2N)$ が開弦でも整合する最初の確認である。

2 箱でも同じ構造を確かめた。C$_N$ の Macdonald 多項式 $P_{2\varepsilon_1}$, $P_{\varepsilon_1+\varepsilon_2}$ を Macdonald 作用素の三角行列から求め、$N = 2, \dots, 9$ の値を $x = t^N$ の有理関数に補間した（$u = 1$ では D$_N$ の 2 箱に一致し、Cherednik の $\mathcal{H}^C_{3,2}(\omega_2)$ に単項式を除いて一致する）。
% [D12][D13][D14]
\begin{compresult}
trefoil の refined Sp$(2N)$ の 2 箱で、
\begin{enumerate}
\item $g^{\mathrm{Sp}}_{(2)} + g^{\mathrm{Sp}}_{(1,1)}$ と $(1+t)(g^{\mathrm{Sp}}_{(2)} - g^{\mathrm{Sp}}_{(1,1)})$ は整数 Laurent 多項式で、SO と同じ全 $d$ の構造（命題~\ref{prop:trace}）をもつ。
\item $E^{\mathrm{Sp}}(t{=}-1) = E^{\SO}(t{=}-1)$（全係数）。
\item hatted の $p_1^2$ 成分 $(g_{(2)} + g_{(1,1)})/z_q$ は Sp$(\ba) = +\,\SO(-\ba)$（[1] は $-$）。$g_R$ そのものは
\begin{align*}
f^{\mathrm{Sp}}_{(2)}(\ba) - f^{\SO}_{(2)}(-\ba) &= \psi^2[\delta], \qquad f^{\mathrm{Sp}}_{(1,1)}(\ba) - f^{\SO}_{(1,1)}(-\ba) = -\psi^2[\delta], \\
\delta &= \mathrm{even}_{\ba}\big(g^{\SO}_{(1)}\big) + 1 = \frac{(\ba^2 - \bq^2)(\ba^2\bt^2 - 1)(\ba^2\bq^2\bt^2 - 1)}{\ba^6\bq^2\bt^4}
\end{align*}
を満たす（$g$ でも同じ。$\psi^2$ は $q, t, \ba^2$ を 2 乗する Adams 作用素）。
\end{enumerate}
\end{compresult}
冪和の成分では $Z_1^{\mathrm{Sp}}(\ba) = -Z_1^{\SO}(-\ba)$, $Z_{11}^{\mathrm{Sp}}(\ba) = Z_{11}^{\SO}(-\ba)$, $Z_2^{\mathrm{Sp}}(\ba) = Z_2^{\SO}(-\ba) + 1$ で、符号は $(-1)^{|\rho|}$、$p_2$ 方向にだけ定数 $1$ のずれがある。$f$ の差 $\psi^2[\delta]$ のうち $\psi^2[\mathrm{even}(g_{(1)})]$ は、[1] の規則が $\ba$ の偶数部分（$c = 1$）の符号だけを変えるのに対し $d = 2$ の Adams 項 $\tfrac12\psi^2[f_{(1)}]$ の偶数部分がそれに従わないことから生じ、$\psi^2[1]$ は $Z_2$ の定数のずれである。彩色の $W_R$ の段階では $W_R^{\mathrm{Sp}}(\ba) = (-1)^{|R|}\,W_R^{\SO}(-\ba)$ が厳密に成り立ち、冪和への変換係数は $p_2$ の定数項（恒等元のトレース）だけが $-2$ 違う。これが $Z_2$ の定数 $1$ で、unknot の [1] のクロスキャップの $\ba$ 偶数部分（SO で $+1$、Sp で $-1$）に等しい。
2 箱の計算は trefoil に限る（$T(2,5)$ は補間に $N \le 13$ が要り、現在の実装では計算量が大きすぎる）。
$E$ が Sp への射影で符号を変えないことは、unrefined の対応「Sp$(N) = \SO(-N)$ かつ $R \to R^{T}$」で $E_{(2)} = -E_{(1,1)}$ を使うと、符号反転と転置が打ち消して $E$ が不変になることと整合する。$p_2$ 方向には [1] の Adams 項 $\psi^2[g_{(1)}]$ が入り、$\psi^2$（$\lambda \to \lambda^2$）は $\ba \to -\ba$ と可換でないため、$f$, $g$ の段階では単純な規則にならない。

\subsection{普遍因子}\label{sec:universal}
% [D9]
\begin{compresult}\label{cr:universal}
refined だけの部分
\[
X_k = \hat g - \hat g\big|_{t=q}(\lambda \to \lambda \qr^k)
\]
が $(q - t)(\lambda^2 - t)$ で割り切れる $k$ を $-3 \le k \le 3$ で調べた。割り切れるのは $k = +1$（$\lambda \qr = \tlam$）だけである。
対象は次のすべてである。
\begin{itemize}
\item [1]：$T(2,3)$, $T(2,5)$
\item 2 箱の $\hat g_{(2)}, \hat g_{(1,1)}$ と和：$T(2,3)$, $T(2,5)$
\item 3 箱の $\hat g_R$：trefoil
\end{itemize}
太字では $(q - t)(\lambda^2 - t) = \bq^2\,(\bt^2 - 1)(\ba^2 - \bq^2)$ で、[1] の $c = 2$ 部分と Hopf 絡み目の $\hat C_{R,S}$ の因子と同じである。
\end{compresult}
$(\lambda^2 - t)$ で割り切れることは、「SO(2) の点 $\lambda^2 = t$ で $\hat g$ が $t$ によらない」ことと同じである。$\lambda$ をずらさないと割り切れない。

\begin{conjecture}\label{conj:structure}
結び目の refined クロスキャップ部分は、どの彩色でも
\[
\hat g_R = \hat g_R\big|_{t=q}(\lambda \to \tlam) + (\bt^2 - 1)(\ba^2 - \bq^2)\cdot(\text{整数係数の Laurent 多項式})
\]
の形をとる。
\end{conjecture}

\begin{remark}
閉弦の凍結した状態と向き付けられた状態の電荷の比は $x^2/Q_0 = \qr^{-1}$ である（AgSc の規約、§\ref{sec:index}）。開弦では、クロスキャップの変数 $\tlam^2$ と被覆の変数 $\ba_{\mathrm{KN}}^2 = \lambda^2w^{-1}$ の比が $w^3$ で、両者は一致しない。
開弦と閉弦の変数の対応（$\lambda^2$ と $Q$ か、$\lambda^{-2}$ と $Q$ か）と開弦の規格化がまだ固定されていないので、この比較の結論は保留する。
\end{remark}

\section{正値性の現状と障害}\label{sec:positivity}

\subsection{符号の解釈}
BPS 数は状態の数そのものではなく、$(-1)^F$ を付けた指数である。
閉弦では $N^{s_L}_\beta = \sum_{s_R}(-1)^{2s_R}(2s_R + 1)N^{s_L,s_R}_\beta$（AgSc §2）、開弦の LMOV 不変量では $\hat N_{\mu,g,\beta} = \sum(-1)^{2J_r}\hat N_{\mu,g,\beta,J_r,S_R}$（KN）と書ける。
KN の正値性は、refined にすると $(-1)^F$ の打ち消しが解けて、スピンで分けた状態の数 $\hat N_{\dots,J_r} \ge 0$ が見えることを指す。

orientifold では、AgSc が指摘するとおり、$\tilde I$ が SU$(2)_L\times$SU$(2)_R$ を U$(1)_1\times$U$(1)_2$ に壊す。そのため、向き付けられない状態は完全なスピン多重項をなさなくてよい。
また Piazzalunga--Uranga \cite{PU14} の解釈では、凍結した状態の数は orientifold の偶奇で重みを付けた差 $n^+ - n^-$ である。
したがってクロスキャップの正値性は、適切な多重項の基底で、各ブロックの係数が unrefined の符号 $\varepsilon_\beta$ と同じ符号にそろうこととして問うべきである（§\ref{sec:criterion}）。
unrefined で符号が混在することは、refined で正値性を期待しない理由にはならない。

なお、SO$(2N)$ の refined CS 不変量そのものは、[1] でも太字の変数で正にならない（trefoil の $R_{(1)}$ は 9 項中 4 項が負）。Cherednik \cite{Cherednik11} も、C--D（部分的に B）型の hyper-polynomial について「最小の重みでも正値性はほとんど無い」と述べている。正値性を問うのは BPS の段階である。

\subsection{多箱での符号}
% [P2][P3]
\begin{compresult}\label{cr:signs}
普遍因子で割った商 $X_1/((\bt^2-1)(\ba^2-\bq^2))$ を太字の変数で書き、$\ba$ の冪ごとのブロックで符号がそろう数を数えた。
\begin{center}
\begin{tabular}{lc}
\toprule
対象 & そろうブロック／全ブロック \\
\midrule
{[1]}（$T(2,3)$, $T(2,5)$） & すべて \\
trefoil の 2 箱の和 $\hat g_{(2)} + \hat g_{(1,1)}$ & 5/10 \\
$T(2,5)$ の 2 箱の和 & 6/18 \\
trefoil の 3 箱 $\hat g_{(3)}, \hat g_{(2,1)}, \hat g_{(1,1,1)}$ & 2/15, 3/14, 2/13 \\
\bottomrule
\end{tabular}
\end{center}
次の取り替えでも、符号はそろわない。
\begin{itemize}
\item 自由エネルギーの右辺の基底：Hall--Littlewood、$P_R(Y; q, t)$、$P_R(Y; t, q)$、$Q_R$、modified Macdonald。$P(Y;t,q)$ では 8/11 まで改善するが、Hopf と固定 $N$ の 3 箱では悪くなるか、整数性が壊れる。
\item 二次量子化の指数：$\tilde Z$ の積表示の指数を、$t$ の冪級数（モードの塔）に展開した係数。混在するブロックの割合は、trefoil の 2 箱で約 13\%（123 中 16、110 中 15）、$T(2,5)$ の 2 箱で約 34\%、trefoil の 3 箱で 38〜44\% である。
\item 多重項の変数：$(q^{1/2}, \qr)$ の対称性の探索。
\end{itemize}
\end{compresult}
命題 \ref{prop:alld} と系 \ref{cor:Efree} により、2 箱の和と 3 箱の方向は $E$ によらない。したがってこの混在は $E$ の不定性によるものではない。凍結した状態 $C$ の多重項の基底がまだ正しくないためと考えられる（予想。正値性そのものが成り立たない可能性は排除していない）。

多箱の $c = 1$ 部分が $\sigma$ 対称でない（$G$ 展開できない）理由は、次の二つに分けられる。
% [P10]
\begin{compresult}\label{cr:multibox-sigma}
\begin{enumerate}
\item $t = q$ では Mariño の $c$ の定義は $z = q^{1/2} - q^{-1/2}$ の偶奇（$c = 1$ が $z$ の偶数冪、$q \leftrightarrow 1/q$ で対称）であり、$\lambda$ の偶奇とは、冪和の成分で $p_2$ を含まない方向（[1]、$p_1^2$、$p_3$, $p_1^3$）で一致し、含む方向（$p_2$、$p_2 p_1$）で逆になる。hatted $\hat g_R = M(q^{-1})^{-1} g_R$ では、$c = 1$ は $\lambda^\beta$ の $\beta \equiv |R| - 1 \pmod 2$ の部分である（[1] は偶、2 箱は奇、3 箱は偶。trefoil の 3 箱の $\hat g_R$ は $U$ 一般で整数 Laurent 多項式で、3 つとも成り立つ）。$p_2$ 方向の $\lambda$ の偶数冪に混ざる $c = 2$ は閉じた形 $-2 z\,\delta|_{t=q}\,(q^{4n}-1)/(q-1) \times$（単項式、$T(2,5)$ では $\lambda^4 q + \lambda^2(q^2+q+1) + q$ を含む）をもつ（$n = 2m+1$、$\delta$ は \S\ref{sec:sp} のもの）。
\item 正しく $c = 1$ を取っても refined では $\sigma$ 対称でない。前の因子、[1] の積や Adams 項の補正、転置では直らない。$\sigma$ 奇部分は $\bq^2(\bt^2 - 1)(V^4 - 1)$（$V^4 = qt$）で割り切れ、$t = q$ で消える。すなわち $\hat g^{c=1} = S + (\bt^2-1)(\ql - \ql^{-1})\,K'$（$S$, $K'$ は $\sigma$ 対称で $G$ 展開できる）。ただし $S$, $K'$ は半整数で、$G$ 展開のブロックの符号は trefoil の $K'$ 側で全ブロックそろう（$\beta$ で交互）が、$T(2,5)$ では崩れる。
\end{enumerate}
\end{compresult}
因子 $(\bt^2 - 1)$ は $c = 2$ の普遍因子（SU(2)$_r$ の捻ったトレース）、$(\ql - \ql^{-1})$ は Lefschetz SU(2)$_\ell$ の Weyl 反射で奇である。多箱の $c = 1$ には「$\tilde I$ の捻りを受け、Lefschetz のスピンが奇」の状態が入り、[1] にはそれが無い。これは \S\ref{sec:EC} の「結び目の固定セクターでは $S_R$ が自明に働かない」の $G$ 展開側の表れと考えられる（予想）。

\subsection{被覆の brane と image の分割}
被覆を brane（$Y$）と image（$Y'$）に分けた生成関数を作ると、次が成り立つ（trefoil, $T(2,5)$）。
\begin{itemize}
\item brane--brane のセクターは、KN の t 枠 $(t^{-1/2} - t^{1/2})\hat f$ で $G$ 展開でき、全 $\beta$ で符号がそろう（KN の正値性）。
\item brane--image のセクター $x = W_{\mathrm{adj}} + \kappa - W_{(1)}^2$ は、$(t^{-1/2} - t^{1/2})x/z_q$ が整数 Laurent で $\sigma$ 対称である。$\lambda^\beta$ のブロックの多くは「そろう」か「$(\bt^2 - 1)\times$ そろう」かのどちらかである。混在するのは、trefoil では $\lambda^{-8}$ と $\lambda^{-2}$、$T(2,5)$ では中央の $\lambda^{-16}$〜$\lambda^{-10}$ の 4 ブロックである。
\end{itemize}
向き付けられた被覆の brane--image のセクターが、クロスキャップと同じ普遍因子 $(\bt^2 - 1)$ をもつ。
brane--image のセクターは、$\tilde I$ で自分に移る状態（$E$ の源）を含むセクターである。

\subsection{$E$ と $C$ を分ける試み}\label{sec:EC}
命題 \ref{prop:EC} により、$E$ と $C$ の分離には指数以外の情報が要る。次の 3 つを試した。
% [P4]
\begin{compresult}\label{cr:EC}
\begin{enumerate}
\item \textbf{状態の部分和}：Hopf と同じく $E = P - 2Q$（$P = (1-t)x$）と書き、(i) $C$ が整数 Laurent、(ii) $t = q$ で $E = -(1+q)(g_{(1)}(q^2, \lambda^2) + 1)$ を課した。trefoil で解が無い（$T(2,5)$ では実行していない）。範囲を外しても、(i)(ii) を満たす $E$ は $P$ の単項式の台に収まらない。
\item \textbf{Adams 型}：$E = F\cdot(\psi^2[g_{(1)}] + 1)$。$F$ は $(1+q)$, $(1+t)$, $(q+t)$, $(1+qt)$ とその単項式倍で、$\psi^2$ で $\lambda^2$ に $(q/t)^{j/2}$（$|j| \le 2$）を掛ける場合も含めた。どれも (i) を満たさない。
\item \textbf{捻ったトレース}：固定セクターの指数を $P = t^{1/2} z_q P_k$ と書き、$P_k$ の多重項（$G$ 展開の項、または単項式）ごとに $\tilde I$ の捻った指標をとる。境界の因子を $X$ として $E = sX\sum_m e_m\tilde b_m$、$|e_m| \le |c_m|$、$e_m \equiv c_m \pmod 2$ と置き、整数計画で解いた。捻りの規則は仮説で、q 平面・t 平面の $\pi$ 回転と、単項式ごとの符号の 3 通りを試した。
  \begin{itemize}
  \item Hopf の $E_{22}, E_{24}$ は、$X = z_q$ と単項式ごとの符号で一意に再現される。
  \item 結び目では、$t = q$ と $t = -1$ の目標値がどちらも $(1 + q)$ を含むので、$X = z_q$ では解が無い。境界の因子は和 $z_q^+$ でなければならない。
  \item $X = z_q^+$ では、trefoil は一意に解けるが、解に規則性は見えない。$T(2,5)$ はどの規則でも解が無い。上限 $|c_m|$ をゆるめると、見えない不変な対が少なくとも 2 組要る。
  \end{itemize}
\end{enumerate}
境界の因子が和になることは、命題 \ref{prop:boundary} で導出した。
\begin{enumerate}
\setcounter{enumi}{3}
\item \textbf{境界因子を固定した整数計画}：境界因子を $z_q^+$ に固定し、モジュライ部分の $G$ 展開の多重項に取っ手の捻り $G\to X$（$X\in\{z_q, z_t, (z_q+z_t)/2, (qt)^{1/2}-(qt)^{-1/2}, z_q^+-z_t^+\}$）を当てても、$T(2,3)$, $T(2,5)$ で解が無い。
\item \textbf{framed unknot と符号の規則}：framed unknot（framing $\mathsf{f}$）では固定セクターの指数が $P(\mathsf{f}) = (\lambda^2-1)(q-1)\sum_{k<\mathsf{f}}(\lambda^2q)^{\mathsf{f}-1-k}t^k$ の積になる。条件 (ii) の形を仮定せず、(i) と「$t = q$ で $C$ が整数 Laurent」だけを課して全候補を数えると、$\mathsf{f} = 1, 2, 3$ のすべてで $E$ は一意に
\[
E(x) = (-1)^x(\lambda^2+1)(q+1)\sum_{k<x}(-1)^k(\lambda^2q)^{x-1-k}t^k\times\text{単項式}
\]
と決まる。$\tilde I$ は $t$ の塔に $(-1)^k$（$e^{i\pi J_2}$）で働く。Hopf $(2,2)$ でも、$t$ と $\lambda$ の指数の偶奇で符号を付ける規則が $E_{22}$ を再現する。しかし Hopf $(2,4)$ と結び目では、この型の規則（8 通りと基準のずれ 4 通り）も、符号と取っ手の捻りを組み合わせた 36 通りの規則も、どれも (i)(ii) を満たさない。
\end{enumerate}
\end{compresult}
$\tilde I = I_X\otimes e^{i\pi J_2}$ であり、指数が見るのは $S_R - J_2$ の組み合わせだけである。したがって $E$ の決定には、固定セクターの状態の $t$ 電荷を $J_2$（$\tilde I$ で符号）と $S_R$（$\tilde I$ で不変）に分ける情報、すなわち KN の $\hat N_{\rho,g,\beta,J_r,S_R}$ に当たる打ち消しの前の状態数が要る。unknot ではすべて $J_2$、結び目では同じ電荷の状態に異なる符号が付く（計算）。KN の仮定（$S_R$ が BPS 状態に自明に働く）のもとでは符号が電荷の偶奇で決まるはずなので、結び目の brane--image のセクターでは $S_R$ が非自明に働くと考えられる（予想）。これが命題 \ref{prop:EC} の不定性の、量子数による言い換えである。

\begin{remark}[refined の強い整数性]
全 $d$ の形で $f = \tfrac12 f^{\cov} + \tfrac12 E/(1+t) + C$ が整数 Laurent であることは、不変な状態の指数 $P_{\mathrm{inv}}$ について $E \equiv P_{\mathrm{inv}} \pmod 2$（係数ごと）と同値である（$\tfrac12[P_{\mathrm{inv}}(1+t) + E(1-t)]/(1-t^2)$ が整数級数 $\Leftrightarrow P_{\mathrm{inv}} + E \equiv 0 \bmod 2$）。これは整数計画の制約 $e_m \equiv c_m \pmod 2$ そのものである。AgSc §4 の unrefined の形（向き付けられた $f_{\square,\square}$ と向き付けられない $f_\square(q^2)$ の mod 2 の一致）は、$t = q$ で $E$ と $C$ が条件 (ii) で混ざった後の影である。その形を素朴に当てた検査 $P \equiv g_{(1)}(q^2,t^2,\lambda^2 w^j) \pmod 2$ は $t = q$ でも合わず、辞書が未確定なので refined についての判定はしていない。
\end{remark}

\subsection{$n = 1$ の $E$：状態の空間の同定}\label{sec:principle}
固定セクターの最低次（brane 1・image 1、$n = 1$）では、$E$ は条件を使わずに状態の空間の同定だけから作れる。$T(2,2m+1)$ の $E$ の台 $N$ を $\lambda^{-2}$ の冪（ブロック $k = 0, \dots, K = 4m-1$）で分け、$V = 1 + t + qt$（trefoil の簡約 superpolynomial）、$S$ を $T(2,2m+1)$ の簡約 superpolynomial、$\psi^2$ を $(q, t)\to(q^2, t^2)$ とする。
% [P17]
\begin{compresult}[符号の規則の破れ]\label{cr:signrule}
下の系列（$k \le 2m$）に $(-1)^{\deg_q}$、上の系列に $(-1)^{\deg_t}$ を付ける符号の規則は、$T(2,3)$ と $T(2,5)$ では条件 (i)（$t = -1$）と (ii)（$t = q$）のもとで $E$ を一意に決めるが、$T(2,7)$ では (i) だけでも解が無い。破れるのは上の系列の最初のブロック $k = 2m+1$ である。
\end{compresult}
% [P18]
\begin{compresult}[状態の空間と $\tilde I$]\label{cr:principle}
$N$ のブロック $n_k$（$\lambda^{-(6+2k)}$ の係数を $r^{-(4+2k)}$ で正規化したもの）は
\[
n_k = \mathrm{Sym}^k V \ (k < 2m),\qquad n_{2m} = \mathrm{Sym}^{2m}V - S\otimes S,\qquad n_{K-j} = -t^{K-2j}\bigl[\mathrm{Sym}^j V + (q-t)\,R_j\bigr]\ (0 \le j \le 2m-2)
\]
であり、$\tilde I$ は $V$ の対合 $\mathrm{diag}(1, 1, -1)$（上の系列では双対）の誘導作用、$S\otimes S$ では因子の入れ替え（トレース $\psi^2 S$）として働く。この同定から作った $E$ は、$T(2,3)$・$T(2,5)$ で (i)+(ii) から決まる一意な $E$ を全単項式で再現し、$T(2,7)$〜$T(2,15)$ では予言として (i)・(ii) をともに満たす。
\end{compresult}
% [P19]
\begin{conjecture}[補正の形]\label{conj:correction}
$R_{m+j'} = t^{m-1-j'}F_{j'}$ で、$F_{j'}$ は $m$ によらず
\[
\sum_{j'\ge 0}F_{j'}z^{j'} = \frac{1 + (2t + t^2 + qt + 2qt^2)z + qt^3z^2}{(1-z)(1-qtz)(1-t^2z)(1-q^2t^2z)}
\]
（状態の数は $\binom{2j'+3}{3}$）。$j' \le 5$（$T(2,7)$〜$T(2,15)$）で全係数が一致する。
\end{conjecture}
% [P20]
\begin{conjecture}[補正への $\tilde I$ の作用]\label{conj:correctionI}
$E$ の補正は $(-1)^{m+j'+1}t^{m-1-j'}\psi^2(X_{j'})$、$X_{j'} = \mathrm{Sym}^{j'}V + t\,\mathrm{Sym}^{j'-1}V$（母関数 $(1+tz)/((1-z)(1-tz)(1-qtz))$、$(j'+1)^2$ 状態）であり、$F_{j'}$ の残りの状態は $\tilde I$ で入れ替わる対である：$F_{j'} = \psi^2 X_{j'} + 2P_{j'}$、$\sum P_{j'}z^{j'} = tz(V' + qt^2z)/((1-z)(1-qtz)(1-t^2z)(1-q^2t^2z))$、$V' = 1 + q + qt$。$j' = 2$ は $T(2,9)$ で (i) の行と台から 32 通りの変種のうち一意に選ばれ、$j' = 3, 4, 5$ は $T(2,11)$, $T(2,13)$, $T(2,15)$ で (i)+(ii) と整合する。
\end{conjecture}
% [P23]
\begin{conjecture}[閉じた式]\label{conj:closedform}
$T(2,2m+1)$ の $n = 1$ の $N$ と $E$ は、結び目によらない普遍的なデータ $V$、$F_{j'}$、$X_{j'}$ と、結び目のデータ $S_m$（簡約 superpolynomial）・$m$ だけで書ける（$K = 4m-1$、$n_k$ は $\nu^{-(6+2k)}r^{-(4+2k)}$ の係数）：
\[
 n_k = \mathrm{Sym}^kV\ (k<2m),\quad n_{2m} = \mathrm{Sym}^{2m}V - S_m^2,\quad n_{K-j} = -t^{K-2j}\bigl[\mathrm{Sym}^jV + (q-t)\,t^{m-1-j'}F_{j'}\bigr]\ (j = m+j' \ge m),
\]
$E$ は対応する $\tilde I$ のトレースである。$T(2,3)$〜$T(2,15)$ で $N$ は全係数で成立し、$E$ は (i)・(ii) を満たす（計算）。
\end{conjecture}
% [P24]
\begin{compresult}\label{cr:plog}
母関数の plethystic log は、$X_{j'}$ では有限（$(1+2t+qt)z - t^2z^2$：$V\oplus t$ の Fock 空間に関係 $t^2z^2$）だが、$F_{j'}$ と上の系列 $\mathrm{Sym}^jV + (q-t)t^{m-1-j'}F_{j'}$ では無限で、上の系列の生成元は $Vz$ と $(q-t)t^{m-1-j'}c_{j'}z^{m+j'}$（$c_0 = 1$、$c_1 = V^2 - t(1+qt)$、$c_2$ から負の係数）である。上の系列は正の BPS 生成元の Fock 空間としては読めない。
\end{compresult}
% [P25]
\begin{compresult}\label{cr:t34}
$n = 1$ の原理のうち、下の系列 $\mathrm{Sym}^kV$、移り目で完全平方を引く構造、$K = rs - 3$、条件 (i)(ii) を満たす $E$ の存在は、2 本撚りでない $T(3,4)$（$\gamma = \tau_+\tau_-^3$、$K = 9$）でも成り立つ。しかし移り目の平方の因子は $\tilde S = (1+qt)(1+t+t^2+qt^2+q^2t^2)$（10 状態）で $T(3,4)$ の簡約 superpolynomial（11 状態）ではなく、上の系列は $-t^{K-2j}[\mathrm{Sym}^jV + (q-t)R_j]$ の形にならない（$j = 0$ で $-t^8(q+t)$）。下の系列を固定した整数計画は (i)+(ii) で解をもつが一意でない（111 単項式中 29 が不定）。移り目の $S\otimes S$ と補正 $F_{j'}$・$X_{j'}$ の構造は 2 本撚りのトーラス結び目に固有である。
\end{compresult}
% [P26]
\begin{proposition}\label{prop:rows}
$x = \nu^{-2}$ とし、$[1]$ の鏡像の $a$ の行を $R_1 = q^{-\delta/2}t^{-\delta/2}x^\delta R(x)$、$R(y) = \sum_j \rho_j y^j$、随伴 (composite $(1;1)$) の行を $R_{\mathrm{adj}} = q^{-2\delta}x^{2\delta}A(x)$、$A(y) = \sum_j \alpha_j y^j$ と書く。$n = 1$ の固定セクター $N = \sum_k n_k x^k$ は
\[
\sum_k n_k x^k = \frac{1 - x^{2\delta}R(qx)^2}{(1-x)(1-tx)(1-qtx)} + \frac{x^{2\delta-1}G(qx)}{1-tx},\qquad A = R^2 + (q-1)(t-1)G
\]
を満たし、$N$ が Laurent であることと $G$ が Laurent であることは同値である。
\end{proposition}
\begin{compresult}\label{cr:rows}
$T(2,3)$〜$T(2,15)$、$T(3,4)$、$T(3,5)$、$T(3,7)$、$T(3,8)$、$T(3,10)$、$T(4,5)$、$T(4,7)$ の 14 のトーラス結び目で $g_0 = 0$、$g_1 = (2\rho_0\rho_1 - \rho_1^2)/q$ が成り立ち、したがって移り目は $n_{2\delta} = \mathrm{Sym}^{2\delta}V - (\rho_0 - \rho_1)^2$ である。平方の因子は常に $[1]$ の最上 2 行 $\rho_0 - \rho_1$ であり、2 本撚りではこれが簡約 superpolynomial に一致し（行が 2 つ）、$T(3,4)$ の $\tilde S$ は 3 行目 $t^3$ を切り捨てたものである。$N$ は随伴と $[1]$ の被覆データだけで決まるので、$T(3,5)$（$\gamma = \tau_+\tau_-\tau_+\tau_-$、$\delta = 4$、$K = 12$）、$T(3,7)$、$T(3,8)$、$T(3,10)$、$T(4,5)$、$T(4,7)$ では $\tilde S$ と $K = rs - 3$ を予言として的中させた。最深ブロックは $n_K = -t^{K-s+2}h_{s-2}(q,t)$（$s = 2, 3, 4$）である。
\end{compresult}
% [P27, P28]
\begin{compresult}\label{cr:threestrand}
3 本撚りの上の系列 $U_j = -n_{K-j}/t^{K-2j}$ は $j \le \delta - 2$ で $(1 + q/t)\mathrm{Sym}^jV + (q-t)C_j$ の形をもち（$C_j$ は結び目に依る）、$j = \delta - 1$ で $t = q$ の完全平方 $-([\lfloor r/2\rfloor]_{t^2}[r']_t)^2/t^{\delta-1}$（$r'$ は $r$ 以下の最大の奇数）が入る（$T(3,4)$、$T(3,5)$、$T(3,7)$、$T(3,8)$、$T(3,10)$）。4 本撚り（$T(4,5)$、$T(4,7)$）では前因子が $h_2(q,t)/t^2$、平方の位置が $j = r - s + 1$ で、平方の因子は $\rho_0(T(s-1, r-s+1))\cdot S(T(2,r'))|_{t=q}$ と読める。$E$ については、移り目の規則 $E_{2\delta} = \mathrm{Sym}^{2\delta}V|_{\varepsilon} - \psi^2\tilde S$ が $T(3,4)$・$T(3,5)$ でも条件 (i)(ii) と両立し、最深の 2 ブロックが一意に決まる（$E_K = -t^8(q-t)$、$+t^{11}(q-t)$）。2 本撚りでは、$E$ を同じ $(R, G)$ の形に分けると $G$ の行の $\tilde I$ は深い行（$j \ge \delta + 2$）で $q \mapsto -q$ の規則、それより上で対の補正 $D_{\delta-i} = t^{a(m,i)}\tilde D_i$（$\tilde D_i$ は $i$ と $m$ の偶奇だけで決まる）である。
\end{compresult}
% [P29]
\begin{compresult}[composite の大きな $N$ の因子化]\label{cr:composite}
命題 \ref{prop:rows} の $A = R^2 + (q-1)(t-1)G$ は随伴に限らない。trefoil の composite $(2;1)$、$(1,1;1)$、$(2;2)$、$(2;1,1)$、$(1,1;1,1)$ の大きな $N$ の値は
$R_{[\alpha;\beta]} = R_\alpha R_\beta + (q-1)(t-1)G_{\alpha\beta}$（最上行は積、差は全行で $(q-1)(t-1)$ で割れる）を満たし、$G_{\alpha\beta}$ の行は $\alpha, \beta$ に応じた次元型の因子 $(1+q)$、$(1+t)$ を含む。$(q-1)(t-1)$ は brane と image の間の環状の弦の refined の測度で、$N$ の整数性は連結相関の LMOV 型の整数性と同値である（読み）。$g_1 = (2\rho_0\rho_1 - \rho_1^2)/q$ の一般化 $\rho_0^\alpha\rho_1^\beta + \rho_1^\alpha\rho_0^\beta - \rho_1^\alpha\rho_1^\beta$ は他の composite では成り立たない。
\end{compresult}
% [P30, P31, P38, P39]
\begin{compresult}[$n = 2$ の行の母関数と $\lambda$ 依存性]\label{cr:n2rows}
trefoil の $n = 2$（brane 2・image 2）の固定セクター $N_\lambda$（$\lambda = (2), (1,1)$）は、$[1]$・2 箱・随伴・composite の行（$a_2, a_4$；$b_4, b_6, b_8$；$c_4, c_6, c_8$；$G$ の行）の多項式（最大 4 次）として明示的に書け、実際の行で $\nu^{-20}$ まで全ブロックが保存した $N_\lambda$ と一致する。
\begin{enumerate}
\item \textbf{$\sigma$ 共変性}：2 箱の行と $G$ の重みの付け替え $\sigma$：$b_j \mapsto qb_j$、$c_j \mapsto c_j/t$、$G_{\alpha\beta} \mapsto w_\alpha w_\beta G_{\alpha\beta}$（$w_{(2)} = q$、$w_{(1,1)} = 1/t$）に対して、$N_{(2)}\circ\sigma - N_{(1,1)}$ は $(q-t)\times$（随伴セクターの記号 $a_2, a_4, g^{\mathrm{adj}}_j$ と背景だけの多項式）になる（$\nu^{-20}$ まで）。すなわち $\lambda$ の違いは「2 箱の行の重みの付け替え」と「真空の 2 状態 $(q-t)\otimes$ 随伴の塔」に分かれる。
\item \textbf{$\nu^{-12}$ の平方}：$\sigma$ 共変な部分は両 $\lambda$ で完全平方 $-(s^3/r)\hat Y_\lambda^2$ で、$\hat Y_{(1,1)} = (1+q)(1+t)(1+qt)/q^3$、$\hat Y_{(2)} = [(1+q)(1+t)(1+qt)^2 + qt(q-t)]/q^4$（10 単項式、正、既約。積の形ではない）。補正は $(q-t)[(t/q^4)a_2^2 + \delta]$、$\delta = [tV'^2a_2^2/q^3 - 2a_2\tilde S + (V'-t)g^{\mathrm{adj}}_6]/q^3$（$V' = 1+q+qt$、$\tilde S = (1+t+qt)/q$。$\lambda = (1,1)$ では $\delta = 0$）。brane/image の行を別の記号にすると、この平方は文字どおり brane の $X$ と image の $X$ のテンソル積 $X^bX^i$（$X$ は $a_2^2, b_4, c_4$ の線形結合）である。
\item \textbf{$\hat Y_\lambda$ は brane 単独の量}：単一の brane の 2 箱のキュムラント（$W_1, W_{(2)}, W_{(1,1)}$ の $\log$ と Adams $f_{(2)} = Z_{(2)} - \tfrac12\psi^2 Z_{(1)}$）を Schur の $\lambda$ に射影した $\Omega^{(0)}_\lambda(\nu) = \sum_\rho \chi_\lambda(\rho)f_\rho/\phi_\rho$ は $\mathrm{fac}\cdot[\Omega'_2\nu^{-2} + \Omega'_4\nu^{-4} + \cdots]$（$\mathrm{fac} = -q^{1/2}t/(t-1)$）で、両 $\lambda$ で厳密に $\Omega'_4 = \hat Y_\lambda$ である。$\hat Y_{(2)}$ の $qt(q-t)$ は Adams の項に由来する。さらに $\lambda = (1,1)$ では補正も $\Omega'_2 = -a_2/q$ の平方で、$N_{12}(1,1) = -(s^3/r)[\Omega_4'^2 + (q-t)(t/q^2)\Omega_2'^2]$ と書ける。
\end{enumerate}
\end{compresult}
% [P40]
\begin{compresult}[随伴の塔の母関数]\label{cr:tower}
$x = \nu^{-2}$、$V = 1+t+qt$、$V' = 1+q+qt$ とする。$N_{(2)}\circ\sigma - N_{(1,1)} = -(s^3/r)(q-t)[R_1(x)^2\kappa_R(x) + G_{\mathrm{adj}}(x)\kappa_G(x) + \beta(x)]$ の核は閉じた形で
\begin{align*}
\kappa_G(x) &= \frac{q^2x^3[\,q^2(V'-t) - t^2x(V-q)\,]}{(tx-q)(t^3x-q)(t^2x-q^2)(tx-q^3)},\qquad
\beta(x) = \frac{-q^2x^4}{(x-q)(tx-q)(t^2x-q)(tx-q^2)(tx-1)},\\
\kappa_R(x) &= \frac{-q^2x^3\,P(x)}{(x-q)(tx-q)(t^2x-q)(t^3x-q)(tx-q^2)(t^2x-q^2)(tx-q^3)(tx-1)}
\end{align*}
である（$P$ は $x$ の 5 次の既約多項式）。2 箱の行と $G$ を 0 にしたパイプラインを $\nu$ の有理式のまま計算して得た（記号ブロックの級数と $\nu^{-20}$ まで一致）。核はすべて「真空の 2 状態 $(q-t)\times$ セルの Fock 核」で、セルは背景の 6 セル（$t, 1/q, t/q, t^2/q, t/q^2$；$t^2/q^2$）に $G_{\mathrm{adj}}$ に特有の $t^3/q, t/q^3$ を加えたものである。$\kappa_G$ の級数の係数は階段型 $t^nT_nU_n/q^{3(n+1)}$、$T_n = \sum_{0\le j\le i\le n}q^it^j$ である。
\end{compresult}
% [P32, P33, P34, P41]
\begin{compresult}[$n = 2$ の $\nu^{-12}$ の $E$：全部入れ替え]\label{cr:n2E}
$n = 2$ の $\nu^{-12}$ の $E$ は両 $\lambda$ で「背景のセルの符号 $+$ 結び目の部分のすべての平方・行の積に入れ替え $\psi^2$」で書け、条件 (i)(ii) を満たす：
\[
E_{12}(1,1) = r^{-13}s^3\bigl[\psi^2\hat Y_{(1,1)} + (q-t)(t/q^2)\,\psi^2\Omega'_2\bigr] - B_{12}(1,1),\qquad
B_{12}(1,1) = -r^{-13}s^3(q-t)qt^2 .
\]
以前の読み（補正は $t\,\mathrm{Sym}^2(1+qt)$）の余分な $(q-t)qt^2$ は背景の $E$（$E_B = -B$）であった。$\lambda = (2)$ では平方 $\psi^2\hat Y_{(2)}$ に加えて補正の各部分（$(t/q^4)a_2^2$、$tV'^2a_2^2/q^6$ は $\psi^2$、交差項 $-2a_2\tilde S$ は 0、$G_{\mathrm{adj}}$ の最上行は $n = 1$ と同じ入れ替え $g_1 = (2\rho_0\rho_1 - \rho_1^2)/q \mapsto -\rho_1^2/q$、その核 $V' - t$ はセルの符号 $(-1)^{\deg_q\deg_t}$ で $(1+q)(1-t)$）が符号まで一意に決まる（核の符号 16 通り中 1 通り）。この $E$ は、補正の単項式ごとに $|e|\le|R|$ を課す以前の族の外にある。ブロックの符号は $n = 1$ の移り目、$\lambda = (2)$ の $\nu^{-10}$、両 $\lambda$ の $\nu^{-12}$ を通じて $(-1)^{k/2+1}$（$k$ は $\nu^{-k}$ の $k$）で統一され、$n = 1$ の $F^\varepsilon$ のセル 1 の符号の反転に当たる。$\kappa_G$ の分子の第 2 項 $V - q$ は同じ規則で $(1-q)(1+t)$ になる（予想）。
\end{compresult}
% [P35, P36, P37, P42]
\begin{compresult}[$\nu^{-14}$ 以降の障害]\label{cr:n2negative}
$\nu^{-8}$〜$\nu^{-12}$ の $E$ を固定した $\nu^{-14}$ の整数計画は、(i)+(ii) でも 113 未知数のうち 2〜3 個しか決めない（入れ替わる対の最小は (i) で 48、(i)+(ii) で 936）。入れ替えのトレースを記号ブロックの単項式ごとに取る構成は、どの記号を brane/image のラベルにしても $E$ にならない（$\lambda$ によらないか、Laurent でない）：$N$ の Laurent 性は内容の異なる単項式の相殺で生じるので、トレースは Schur 基底の brane の対象 $\hat Y_\lambda = \Omega'_4$ でしか定義できない。純粋な brane$\times$image の積は混合キュムラントで相殺し（階数 0）、平方 $\hat Y_\lambda^2$ は連結の核（composite の次元の補正、$c_1$、$\kappa$）が担う。$\sigma$ 共変な部分は $\nu^{-14}$ で $\hat Y_\lambda$ で割れず、$\Omega^{(0)}$ の準位の積の級数では書けない。
\end{compresult}
% [P21]
\begin{proposition}\label{prop:pairblind}
$(q-t)$ の補正の対 $(q\mu, -t\mu)$ の $\tilde I$ の符号は条件 (ii) に現れない（$t = q$ で消える）。したがって補正の符号は (i) だけからは決まらず（$n = 1$ の $j' = 2$ で 15 中 2 個、$n = 2$ の $\lambda = (2)$ で 17 中 3 個）、予想 \ref{conj:correctionI} のような状態の空間の構造が要る。
\end{proposition}
したがって $n = 1$ では、$E$ と $C$ の分離（命題 \ref{prop:EC}）は状態の空間の同定で解消する。$n = 2$（brane 2・image 2）では、$N_\lambda$ が Wilson loop の無い背景（6 セルの Fock 空間、閉じた形）と結び目の部分に分かれ、$\nu^{-12}$ までは「背景のセルの符号 $+$ 平方と行の積に入れ替え」で両 $\lambda$ の $E$ が符号まで決まる（計算結果 \ref{cr:n2E}）。$\nu^{-14}$ 以降は、$\sigma$ 共変な部分を brane の対象の平方の和として書く原理が無く、$E$ は未決定である（計算結果 \ref{cr:n2negative}、\S\ref{sec:discussion}）。

\subsection{composite の DAHA の非正値性}
% [P5]
固定セクターの状態は、被覆の随伴 $(1;1)$ の値 $R_{\mathrm{adj}}$ から来る。$R_{\mathrm{adj}}$ がホモロジーの Poincaré 多項式（係数が正）として読めれば、打ち消しの前の状態の数が得られるはずである。
\begin{compresult}
大きな $N$ の $R_{\mathrm{adj}}$（GL の DAHA）は、KN の太字変換を鏡像で当てても、係数の符号が混在する（trefoil は 19 項中 7 項、$T(2,5)$ は 56 項中 25 項が負）。
単項式の偶奇によるどの符号の付け方（64 通り）でも、係数はそろわない。
\end{compresult}
これは Cherednik--Elliot \cite{CE15} の序論の記述「composite の DAHA 超多項式は正値性をもたない」と合う。
環状の HOMFLY-PT ホモロジーとの関係は、同論文でも予想されていない。

\subsection{まとめ}
% [P6]
\begin{itemize}
\item [1]（トーラス結び目）では判定式（予想 \ref{conj:criterion}）が成り立ち、符号がそろう。
\item 多箱では、普遍因子（予想 \ref{conj:structure}）は成り立つが、その商の符号はそろわない。
\item $n = 1$ の $E$ は状態の空間の同定（計算結果 \ref{cr:principle}、予想 \ref{conj:correction}・\ref{conj:correctionI}）から条件なしに作れ、$T(2,3)$〜$T(2,15)$ で (i)・(ii) を満たす。
\item $n = 2$（trefoil）では、$N_\lambda$ の $\lambda$ 依存性が「2 箱の行の重みの付け替え」と「真空 $(q-t)\otimes$ 随伴の塔」に分かれ（計算結果 \ref{cr:n2rows}・\ref{cr:tower}）、$\nu^{-12}$ の平方の因子 $\hat Y_\lambda$ は brane 単独の $\lambda$ 射影キュムラントの第 2 準位である。$\nu^{-12}$ の $E$ は両 $\lambda$ で全部入れ替えとして決まる（計算結果 \ref{cr:n2E}）が、$\nu^{-14}$ 以降は未決定である（計算結果 \ref{cr:n2negative}）。
\item 多箱の正値性は、凍結した状態の多重項の基底が決まるまで判定できない。$E$ と $C$ の分離が決まるまでも同様である。$E$ と $C$ の分離は指数からは決まらず（命題 \ref{prop:EC}）、composite の DAHA の値からも状態の数は得られなかった（計算）。打ち消しの前の固定セクターの状態の空間を、幾何から求める必要がある \open。
\item 多箱の $c = 1$ の $\sigma$ 非対称性は、unrefined の段階での $c$ の偶奇の取り違えと、refined に固有の $\sigma$ 奇部分（$(\bt^2-1)(V^2-V^{-2})$ で割り切れる）に分かれる（計算結果 \ref{cr:multibox-sigma}）。後者の起源は、結び目の固定セクターで $S_R$ が非自明に働くことにあると考えられる（予想）。
\end{itemize}

\section{議論}\label{sec:discussion}

\subsection{未解決の問題}
\begin{enumerate}
\item \textbf{凍結した状態の多重項}：多箱の $c = 1$ 部分は $\sigma$ 対称でなく、$G$ では展開できない。$G = T_+T_-$ の因子 $T_\pm$ を使っても、$\sigma$ 対称でない部分は展開できるようにならない（§\ref{sec:index} の注意）。
向き付けられない面の refined の取っ手の因子を幾何から決める必要がある。
\item \textbf{$E$ と $C$ の分離}：指数からは決まらない（命題 \ref{prop:EC}）。Hopf の $(4,4)$ と結び目の 2 箱で未決定である。
固定セクター（被覆の環状の弦、商では Möbius の帯）の Hilbert 空間への $\tilde I$ の作用を、指数ではなく状態の空間の上で構成する必要がある。
\item \textbf{普遍因子の意味}：$(\ba^2 - \bq^2)$ は SO(2) の点で、$(\bt^2 - 1)$ は $t = q$ と偶奇で説明できる。ただし、$\nu \to \tnu$ のずれの向きが一意に決まる理由と、閉弦の Kähler 変数の三重化との関係は未解明である。
\item \textbf{他の結び目}：非トーラス結び目では、refined SO の値に定義が無い（refined CS の $S, T$ が使えない）。
固定 $N$ では、SO(4) $=$ SU(2)$^2$、SO(6) $=$ SU(4) の同一視で A 型の入力（KNTZ \cite{KNTZ19} の長方形彩色の式など）を使える。二重ツイスト結び目 7 個で、$N = 2$ の SO 側（$|R| \le 3$）の hatted の整数性と一様な規格化が成り立つことを確かめた。
大きな $N$ の値には、D 型の DAHA の iterated torus 結び目への拡張が要る。
\item \textbf{Sp と B 型}：本論文は SO$(2N)$（D 型）に限る。閉弦の B 型（AM）は $c = 1$ だけで、分母に $(t^{n/2} + t^{-n/2})$ をもつ。命題 \ref{prop:alld} の分類で B 型をどう読むかは調べていない。
\end{enumerate}

\subsection{KN との関係}
被覆（向き付けられた部分）の brane--brane のセクターは、KN の $G$ 展開と正値性をそのまま満たす。
破れは、$\tilde I$ に関わるセクター（brane--image とクロスキャップ）に限られる。
したがって、refined の Kauffman 不変量の正値性の問題は、KN の枠組みに orientifold の偶奇の構造を加える問題に帰着する（予想）。


\appendix
\section{コード・データ・テスト}\label{app:code}
\begin{itemize}
\item 実装：\texttt{src/rcsknots/refined\_kauffman/}（\texttt{kauffman\_*.py}、\texttt{koornwinder.py}、\texttt{daha\_dn.py}、\texttt{daha\_cn.py}（C$_N$ 型）、\texttt{daha\_fast.py}、\texttt{flint\_rf*.py}）と \texttt{scripts/refined\_kauffman/rk\_*.py}（Hilbert scheme のセル分解は \texttt{rk\_b\_planB\_cells.py}）。
\item 自動テスト：\texttt{tests/refined\_kauffman/}。保存済みの大きな計算結果を判定し直すテストには \texttt{results} の印を付け、既定の実行から外している。
\item 主張とテストの対応は \texttt{docs/refined\_kauffman/paper/claims.md} に置いた。
\item 以下の式は、保存済みのデータとコードから \texttt{scripts/rk\_report\_tex.py} で生成した（手で写していない）。
\end{itemize}

\section{トーラス結び目の表現 [1]}\label{app:torus}
$g^{\mathrm{true}}_{\yone}$（太字の変数）、$c = 1$ 部分の $G$ 展開、$c = 2$ 部分の $\Pi_m$。
$m = 1$（$T(2,3)$、11 項）：$g^{\mathrm{true}}_{[1]}$ の $\ba$ の冪ごとの係数
\begin{align*}
[\ba^{-6}]\colon\ & \bt^{-4} \\
[\ba^{-5}]\colon\ & \bq \bt^{-2} -\bq \bt^{-4} \\
[\ba^{-4}]\colon\ & -\bq^{2} \bt^{-2} -\bt^{-2} -\bq^{-2} \bt^{-4} \\
[\ba^{-3}]\colon\ & -\bq^{-1} \bt^{-2} +\bq^{-1} \bt^{-4} \\
[\ba^{-2}]\colon\ & \bq^{2} +\bt^{-2} +\bq^{-2} \bt^{-2}
\end{align*}

$m = 1$：$c = 1$ 部分の $G$ 展開（$\nu^\beta$ ごと、$W = (q/t)^{1/4}$）
\begin{align*}
[\nu^{-6}]\colon\ & W^{-8} \\
[\nu^{-4}]\colon\ & -G W^{-6} -2 W^{-4} -W^{-8} \\
[\nu^{-2}]\colon\ & G W^{-2} +1 +2 W^{-4}
\end{align*}

\begin{align*}
\Pi_{1} &= 1
\end{align*}

$m = 2$（$T(2,5)$、27 項）：$g^{\mathrm{true}}_{[1]}$ の $\ba$ の冪ごとの係数
\begin{align*}
[\ba^{-10}]\colon\ & \bt^{-8} \\
[\ba^{-9}]\colon\ & \bq \bt^{-6} -\bq \bt^{-8} \\
[\ba^{-8}]\colon\ & \bq^{2} \bt^{-4} -\bq^{2} \bt^{-6} +\bq^{-2} \bt^{-6} -\bq^{-2} \bt^{-8} \\
[\ba^{-7}]\colon\ & \bq^{3} \bt^{-2} -\bq^{3} \bt^{-4} +\bq^{-1} \bt^{-4} -2 \bq^{-1} \bt^{-6} +\bq^{-1} \bt^{-8} \\
[\ba^{-6}]\colon\ & -\bq^{4} \bt^{-2} -\bq^{2} \bt^{-2} -2 \bt^{-4} +\bt^{-6} -\bq^{-2} \bt^{-4} -\bq^{-4} \bt^{-6} \\
[\ba^{-5}]\colon\ & -\bq \bt^{-2} +\bq \bt^{-4} -\bq^{-3} \bt^{-4} +\bq^{-3} \bt^{-6} \\
[\ba^{-4}]\colon\ & \bq^{4} +\bq^{2} \bt^{-2} +\bt^{-2} +\bq^{-2} \bt^{-4} +\bq^{-4} \bt^{-4}
\end{align*}

$m = 2$：$c = 1$ 部分の $G$ 展開（$\nu^\beta$ ごと、$W = (q/t)^{1/4}$）
\begin{align*}
[\nu^{-10}]\colon\ & W^{-16} \\
[\nu^{-8}]\colon\ & G W^{-10} -G W^{-14} +W^{-8} -W^{-16} \\
[\nu^{-6}]\colon\ & -G^{2} W^{-8} -3 G W^{-6} -2 G W^{-10} -2 W^{-4} -3 W^{-8} \\
[\nu^{-4}]\colon\ & G^{2} W^{-4} +2 G W^{-2} +3 G W^{-6} +1 +2 W^{-4} +2 W^{-8}
\end{align*}

\begin{align*}
\Pi_{2} &= \ba^{2} \bq^{4} \bt^{4} +\ba^{2} \bt^{2} +\bq^{2}
\end{align*}

$m = 3$（$T(2,7)$、49 項）：$g^{\mathrm{true}}_{[1]}$ の $\ba$ の冪ごとの係数
\begin{align*}
[\ba^{-14}]\colon\ & \bt^{-12} \\
[\ba^{-13}]\colon\ & \bq \bt^{-10} -\bq \bt^{-12} \\
[\ba^{-12}]\colon\ & \bq^{2} \bt^{-8} -\bq^{2} \bt^{-10} +\bq^{-2} \bt^{-10} -\bq^{-2} \bt^{-12} \\
[\ba^{-11}]\colon\ & \bq^{3} \bt^{-6} -\bq^{3} \bt^{-8} +\bq^{-1} \bt^{-8} -2 \bq^{-1} \bt^{-10} +\bq^{-1} \bt^{-12} \\
[\ba^{-10}]\colon\ & \bq^{4} \bt^{-4} -\bq^{4} \bt^{-6} +\bt^{-6} -2 \bt^{-8} +\bt^{-10} +\bq^{-4} \bt^{-8} \\
& \quad -\bq^{-4} \bt^{-10} \\
[\ba^{-9}]\colon\ & \bq^{5} \bt^{-2} -\bq^{5} \bt^{-4} +\bq \bt^{-4} -2 \bq \bt^{-6} +\bq \bt^{-8} +\bq^{-3} \bt^{-6} \\
& \quad -2 \bq^{-3} \bt^{-8} +\bq^{-3} \bt^{-10} \\
[\ba^{-8}]\colon\ & -\bq^{6} \bt^{-2} -\bq^{4} \bt^{-2} -2 \bq^{2} \bt^{-4} +\bq^{2} \bt^{-6} -\bt^{-4} -2 \bq^{-2} \bt^{-6} \\
& \quad +\bq^{-2} \bt^{-8} -\bq^{-4} \bt^{-6} -\bq^{-6} \bt^{-8} \\
[\ba^{-7}]\colon\ & -\bq^{3} \bt^{-2} +\bq^{3} \bt^{-4} -\bq^{-1} \bt^{-4} +\bq^{-1} \bt^{-6} -\bq^{-5} \bt^{-6} +\bq^{-5} \bt^{-8} \\
[\ba^{-6}]\colon\ & \bq^{6} +\bq^{4} \bt^{-2} +\bq^{2} \bt^{-2} +\bt^{-4} +\bq^{-2} \bt^{-4} +\bq^{-4} \bt^{-6} \\
& \quad +\bq^{-6} \bt^{-6}
\end{align*}

$m = 3$：$c = 1$ 部分の $G$ 展開（$\nu^\beta$ ごと、$W = (q/t)^{1/4}$）
\begin{align*}
[\nu^{-14}]\colon\ & W^{-24} \\
[\nu^{-12}]\colon\ & G W^{-18} -G W^{-22} +W^{-16} -W^{-24} \\
[\nu^{-10}]\colon\ & G^{2} W^{-12} -G^{2} W^{-16} +2 G W^{-10} -2 G W^{-18} +W^{-8} -W^{-16} \\
[\nu^{-8}]\colon\ & -G^{3} W^{-10} -4 G^{2} W^{-8} -3 G^{2} W^{-12} -5 G W^{-6} -7 G W^{-10} -2 G W^{-14} \\
& \quad -2 W^{-4} -3 W^{-8} -2 W^{-12} \\
[\nu^{-6}]\colon\ & G^{3} W^{-6} +3 G^{2} W^{-4} +4 G^{2} W^{-8} +3 G W^{-2} +6 G W^{-6} +5 G W^{-10} \\
& \quad +1 +2 W^{-4} +2 W^{-8} +2 W^{-12}
\end{align*}

\begin{align*}
\Pi_{3} &= \ba^{4} \bq^{8} \bt^{8} +\ba^{4} \bq^{4} \bt^{6} +\ba^{4} \bt^{4} +\ba^{2} \bq^{6} \bt^{4} +\ba^{2} \bq^{2} \bt^{2} +\bq^{4}
\end{align*}


\section{DAHA：一般の $N$ のreduced 不変量}\label{app:daha}
$\Rred_R(t^N) = T_R^{-6}\times$ 鏡像$(Z)/K_R(t^\rho)$（trefoil、一様な規格化）。$q, t$ の冪は半整数のこともある。
$R_{(1)}(U)$（$U^j$ の係数）：
\begin{align*}
[U^{-5}]\colon\ & -q^{-2} t^{5} +q^{-2} t^{4} \\
[U^{-4}]\colon\ & -q^{-1} t^{4} +q^{-2} t^{4} -q^{-2} t^{3} \\
[U^{-3}]\colon\ & q^{-1} t^{3} -q^{-1} t^{2} \\
[U^{-2}]\colon\ & t^{2} +q^{-1} t
\end{align*}

$R_{(2)}(U)$（$U^j$ の係数）：
\begin{align*}
[U^{-10}]\colon\ & q^{-9} t^{10} -q^{-9} t^{9} -q^{-10} t^{9} +q^{-10} t^{8} \\
[U^{-9}]\colon\ & q^{-7} t^{9} -q^{-7} t^{8} -q^{-9} t^{9} +2 q^{-9} t^{8} -q^{-9} t^{7} +q^{-10} t^{8} \\
& \quad -q^{-10} t^{7} \\
[U^{-8}]\colon\ & q^{-5} t^{8} -q^{-6} t^{8} +q^{-6} t^{7} -2 q^{-7} t^{8} +3 q^{-7} t^{7} -q^{-7} t^{6} \\
& \quad +q^{-8} t^{7} -q^{-8} t^{6} -q^{-9} t^{7} +q^{-9} t^{6} \\
[U^{-7}]\colon\ & -q^{-4} t^{7} +q^{-4} t^{6} -2 q^{-5} t^{7} +2 q^{-5} t^{6} -q^{-6} t^{6} +q^{-6} t^{5} \\
& \quad +q^{-7} t^{7} -3 q^{-7} t^{6} +2 q^{-7} t^{5} -q^{-8} t^{6} +q^{-8} t^{5} \\
[U^{-6}]\colon\ & -q^{-2} t^{6} -q^{-3} t^{6} +q^{-4} t^{6} -2 q^{-4} t^{5} +q^{-5} t^{6} -3 q^{-5} t^{5} \\
& \quad +q^{-5} t^{4} +q^{-6} t^{6} -q^{-6} t^{4} +q^{-7} t^{5} -q^{-7} t^{4} \\
[U^{-5}]\colon\ & q^{-2} t^{5} -q^{-2} t^{4} +q^{-3} t^{5} -q^{-3} t^{4} +q^{-4} t^{4} -q^{-4} t^{3} \\
& \quad +q^{-5} t^{4} -q^{-5} t^{3} \\
[U^{-4}]\colon\ & t^{4} +q^{-2} t^{3} +q^{-3} t^{3} +q^{-4} t^{2}
\end{align*}

$R_{(1,1)}(U)$（$U^j$ の係数）：
\begin{align*}
[U^{-10}]\colon\ & q^{-4} t^{16} -q^{-4} t^{15} -q^{-5} t^{15} +q^{-5} t^{14} \\
[U^{-9}]\colon\ & q^{-3} t^{15} -q^{-3} t^{13} -q^{-4} t^{14} +q^{-4} t^{13} +q^{-4} t^{12} -q^{-4} t^{11} \\
[U^{-8}]\colon\ & q^{-2} t^{13} -2 q^{-3} t^{13} +2 q^{-3} t^{11} +q^{-4} t^{12} -q^{-4} t^{11} -q^{-4} t^{10} \\
& \quad +q^{-4} t^{9} \\
[U^{-7}]\colon\ & -q^{-2} t^{12} -q^{-2} t^{11} +q^{-2} t^{10} +q^{-2} t^{9} +q^{-3} t^{11} -q^{-3} t^{10} \\
& \quad -2 q^{-3} t^{9} +q^{-3} t^{8} +q^{-3} t^{7} \\
[U^{-6}]\colon\ & -q^{-1} t^{10} -q^{-1} t^{9} +q^{-2} t^{10} +q^{-2} t^{9} -3 q^{-2} t^{8} -q^{-2} t^{7} \\
& \quad +q^{-3} t^{9} +q^{-3} t^{7} -q^{-3} t^{6} -q^{-3} t^{5} \\
[U^{-5}]\colon\ & q^{-1} t^{8} -q^{-1} t^{6} +q^{-2} t^{6} -q^{-2} t^{4} \\
[U^{-4}]\colon\ & t^{6} +q^{-1} t^{5} +q^{-1} t^{4} +q^{-2} t^{2}
\end{align*}

被覆の随伴 $(1;1)$ の $R_{\mathrm{adj}}$（$A^{1/2} = \nu$ の冪の係数）：
\begin{align*}
[A^{-5}]\colon\ & q^{-4} t^{5} -q^{-4} t^{4} -q^{-5} t^{5} +q^{-5} t^{4} \\
[A^{-9/2}]\colon\ & 0 \\
[A^{-4}]\colon\ & q^{-2} t^{4} -2 q^{-3} t^{4} +2 q^{-3} t^{3} +q^{-4} t^{4} -2 q^{-4} t^{3} +q^{-4} t^{2} \\
[A^{-7/2}]\colon\ & 0 \\
[A^{-3}]\colon\ & -2 q^{-1} t^{3} +q^{-2} t^{3} -3 q^{-2} t^{2} +q^{-3} t^{3} +q^{-3} t^{2} -2 q^{-3} t \\
[A^{-5/2}]\colon\ & 0 \\
[A^{-2}]\colon\ & t^{2} +2 q^{-1} t +q^{-2}
\end{align*}


\section{大きな $N$ の 2 箱のクロスキャップ部分}\label{app:largeN}
trefoil の $\hat g_{\ytwo} + \hat g_{\yoneone}$、$(1+t)(\hat g_{\ytwo} - \hat g_{\yoneone})$、普遍因子で割った商（太字の変数）。
整数性（計算結果 \ref{cr:largeN}）と普遍因子（§\ref{sec:universal}）の確認用で、$\ba$ の冪ごとの符号はそろわない（計算結果 \ref{cr:signs}）。正値性が見えるセクターと見えないセクターの比較は\ref{app:cover}。
$\hat g_{(2)} + \hat g_{(1,1)}$（整数 Laurent、262 項）：
\begin{align*}
[\ba^{-12}]\colon\ & -\bq^{11} \bt^{-9} -\bq^{9} \bt^{-9} -\bq^{7} \bt^{-9} -\bq^{5} \bt^{-9} -\bq^{3} \bt^{-9} -\bq \bt^{-9} \\
& \quad +\bq^{-1} \bt^{-9} +\bq^{-3} \bt^{-11} +\bq^{-5} \bt^{-13} +\bq^{-7} \bt^{-15} +\bq^{-9} \bt^{-17} +\bq^{-11} \bt^{-19} \\
[\ba^{-11}]\colon\ & -\bq^{12} \bt^{-7} +\bq^{12} \bt^{-9} -2 \bq^{10} \bt^{-7} +\bq^{10} \bt^{-9} -2 \bq^{8} \bt^{-7} -2 \bq^{6} \bt^{-7} \\
& \quad -\bq^{6} \bt^{-9} -2 \bq^{4} \bt^{-7} -2 \bq^{4} \bt^{-9} -2 \bq^{2} \bt^{-7} -2 \bq^{2} \bt^{-9} -\bq^{2} \bt^{-11} \\
& \quad -4 \bt^{-9} -2 \bt^{-11} -4 \bq^{-2} \bt^{-11} -\bq^{-2} \bt^{-13} -4 \bq^{-4} \bt^{-13} -3 \bq^{-6} \bt^{-15} \\
& \quad -2 \bq^{-8} \bt^{-17} -\bq^{-10} \bt^{-19} \\
[\ba^{-10}]\colon\ & \bq^{13} \bt^{-7} -\bq^{11} \bt^{-5} +3 \bq^{11} \bt^{-7} -2 \bq^{9} \bt^{-5} +3 \bq^{9} \bt^{-7} +2 \bq^{9} \bt^{-9} \\
& \quad -3 \bq^{7} \bt^{-5} +3 \bq^{7} \bt^{-7} +3 \bq^{7} \bt^{-9} -4 \bq^{5} \bt^{-5} +2 \bq^{5} \bt^{-7} +5 \bq^{5} \bt^{-9} \\
& \quad -4 \bq^{3} \bt^{-5} +\bq^{3} \bt^{-7} +5 \bq^{3} \bt^{-9} +\bq^{3} \bt^{-11} -2 \bq \bt^{-5} -5 \bq \bt^{-7} \\
& \quad +5 \bq \bt^{-9} +3 \bq \bt^{-11} -4 \bq^{-1} \bt^{-7} -2 \bq^{-1} \bt^{-9} +4 \bq^{-1} \bt^{-11} +\bq^{-1} \bt^{-13} \\
& \quad -5 \bq^{-3} \bt^{-9} -2 \bq^{-3} \bt^{-11} +4 \bq^{-3} \bt^{-13} -4 \bq^{-5} \bt^{-11} -\bq^{-5} \bt^{-13} +2 \bq^{-5} \bt^{-15} \\
& \quad -4 \bq^{-7} \bt^{-13} +\bq^{-7} \bt^{-17} -3 \bq^{-9} \bt^{-15} -2 \bq^{-11} \bt^{-17} -\bq^{-13} \bt^{-19} \\
[\ba^{-9}]\colon\ & 2 \bq^{12} \bt^{-5} -\bq^{12} \bt^{-7} +4 \bq^{10} \bt^{-5} -\bq^{10} \bt^{-9} +5 \bq^{8} \bt^{-5} +2 \bq^{8} \bt^{-7} \\
& \quad -\bq^{8} \bt^{-9} +6 \bq^{6} \bt^{-5} +5 \bq^{6} \bt^{-7} -\bq^{6} \bt^{-9} +6 \bq^{4} \bt^{-5} +9 \bq^{4} \bt^{-7} \\
& \quad +4 \bq^{2} \bt^{-5} +13 \bq^{2} \bt^{-7} +4 \bq^{2} \bt^{-9} -\bq^{2} \bt^{-11} +12 \bt^{-7} +10 \bt^{-9} \\
& \quad +15 \bq^{-2} \bt^{-9} +6 \bq^{-2} \bt^{-11} -\bq^{-2} \bt^{-13} +13 \bq^{-4} \bt^{-11} +2 \bq^{-4} \bt^{-13} +10 \bq^{-6} \bt^{-13} \\
& \quad +6 \bq^{-8} \bt^{-15} +3 \bq^{-10} \bt^{-17} +\bq^{-12} \bt^{-19} \\
[\ba^{-8}]\colon\ & -\bq^{13} \bt^{-5} +\bq^{11} \bt^{-3} -2 \bq^{11} \bt^{-5} -\bq^{11} \bt^{-7} +2 \bq^{9} \bt^{-3} -\bq^{9} \bt^{-5} \\
& \quad -4 \bq^{9} \bt^{-7} +4 \bq^{7} \bt^{-3} -6 \bq^{7} \bt^{-7} -\bq^{7} \bt^{-9} +5 \bq^{5} \bt^{-3} +4 \bq^{5} \bt^{-5} \\
& \quad -10 \bq^{5} \bt^{-7} -2 \bq^{5} \bt^{-9} +4 \bq^{3} \bt^{-3} +10 \bq^{3} \bt^{-5} -11 \bq^{3} \bt^{-7} -6 \bq^{3} \bt^{-9} \\
& \quad +17 \bq \bt^{-5} -7 \bq \bt^{-7} -10 \bq \bt^{-9} -\bq \bt^{-11} +5 \bq^{-1} \bt^{-5} +14 \bq^{-1} \bt^{-7} \\
& \quad -13 \bq^{-1} \bt^{-9} -5 \bq^{-1} \bt^{-11} +8 \bq^{-3} \bt^{-7} +7 \bq^{-3} \bt^{-9} -11 \bq^{-3} \bt^{-11} -\bq^{-3} \bt^{-13} \\
& \quad +7 \bq^{-5} \bt^{-9} +3 \bq^{-5} \bt^{-11} -7 \bq^{-5} \bt^{-13} +5 \bq^{-7} \bt^{-11} +\bq^{-7} \bt^{-13} -3 \bq^{-7} \bt^{-15} \\
& \quad +4 \bq^{-9} \bt^{-13} -\bq^{-9} \bt^{-17} +2 \bq^{-11} \bt^{-15} +\bq^{-13} \bt^{-17} \\
[\ba^{-7}]\colon\ & -\bq^{12} \bt^{-3} -2 \bq^{10} \bt^{-3} -2 \bq^{10} \bt^{-5} +\bq^{10} \bt^{-7} -4 \bq^{8} \bt^{-3} -4 \bq^{8} \bt^{-5} \\
& \quad +\bq^{8} \bt^{-7} -6 \bq^{6} \bt^{-3} -8 \bq^{6} \bt^{-5} +\bq^{6} \bt^{-7} -6 \bq^{4} \bt^{-3} -15 \bq^{4} \bt^{-5} \\
& \quad -\bq^{4} \bt^{-7} +\bq^{4} \bt^{-9} -2 \bq^{2} \bt^{-3} -22 \bq^{2} \bt^{-5} -7 \bq^{2} \bt^{-7} +2 \bq^{2} \bt^{-9} \\
& \quad -11 \bt^{-5} -22 \bt^{-7} +\bt^{-11} -19 \bq^{-2} \bt^{-7} -12 \bq^{-2} \bt^{-9} +2 \bq^{-2} \bt^{-11} \\
& \quad -17 \bq^{-4} \bt^{-9} -5 \bq^{-4} \bt^{-11} +\bq^{-4} \bt^{-13} -12 \bq^{-6} \bt^{-11} -\bq^{-6} \bt^{-13} -7 \bq^{-8} \bt^{-13} \\
& \quad -3 \bq^{-10} \bt^{-15} -\bq^{-12} \bt^{-17} \\
[\ba^{-6}]\colon\ & \bq^{11} \bt^{-5} -\bq^{9} \bt^{-3} +2 \bq^{9} \bt^{-5} -\bq^{7} \bt^{-1} -2 \bq^{7} \bt^{-3} +3 \bq^{7} \bt^{-5} \\
& \quad +\bq^{7} \bt^{-7} -\bq^{5} \bt^{-1} -6 \bq^{5} \bt^{-3} +5 \bq^{5} \bt^{-5} +3 \bq^{5} \bt^{-7} -11 \bq^{3} \bt^{-3} \\
& \quad +4 \bq^{3} \bt^{-5} +8 \bq^{3} \bt^{-7} -7 \bq \bt^{-3} -7 \bq \bt^{-5} +13 \bq \bt^{-7} +2 \bq \bt^{-9} \\
& \quad -19 \bq^{-1} \bt^{-5} +11 \bq^{-1} \bt^{-7} +7 \bq^{-1} \bt^{-9} -4 \bq^{-3} \bt^{-5} -11 \bq^{-3} \bt^{-7} +12 \bq^{-3} \bt^{-9} \\
& \quad +2 \bq^{-3} \bt^{-11} -5 \bq^{-5} \bt^{-7} -4 \bq^{-5} \bt^{-9} +8 \bq^{-5} \bt^{-11} -3 \bq^{-7} \bt^{-9} -\bq^{-7} \bt^{-11} \\
& \quad +3 \bq^{-7} \bt^{-13} -2 \bq^{-9} \bt^{-11} +\bq^{-9} \bt^{-15} -\bq^{-11} \bt^{-13} \\
[\ba^{-5}]\colon\ & \bq^{10} \bt^{-3} +\bq^{8} \bt^{-1} +2 \bq^{8} \bt^{-3} +2 \bq^{6} \bt^{-1} +5 \bq^{6} \bt^{-3} +2 \bq^{4} \bt^{-1} \\
& \quad +10 \bq^{4} \bt^{-3} +2 \bq^{4} \bt^{-5} -\bq^{4} \bt^{-7} +12 \bq^{2} \bt^{-3} +9 \bq^{2} \bt^{-5} -2 \bq^{2} \bt^{-7} \\
& \quad +3 \bt^{-3} +21 \bt^{-5} -\bt^{-7} -\bt^{-9} +9 \bq^{-2} \bt^{-5} +12 \bq^{-2} \bt^{-7} \\
& \quad -2 \bq^{-2} \bt^{-9} +10 \bq^{-4} \bt^{-7} +4 \bq^{-4} \bt^{-9} -\bq^{-4} \bt^{-11} +6 \bq^{-6} \bt^{-9} +\bq^{-6} \bt^{-11} \\
& \quad +3 \bq^{-8} \bt^{-11} +\bq^{-10} \bt^{-13} \\
[\ba^{-4}]\colon\ & \bq^{5} \bt^{-1} -\bq^{5} \bt^{-5} +\bq^{3} \bt^{-1} +2 \bq^{3} \bt^{-3} -3 \bq^{3} \bt^{-5} +7 \bq \bt^{-3} \\
& \quad -6 \bq \bt^{-5} -\bq \bt^{-7} +3 \bq^{-1} \bt^{-3} +\bq^{-1} \bt^{-5} -4 \bq^{-1} \bt^{-7} +6 \bq^{-3} \bt^{-5} \\
& \quad -5 \bq^{-3} \bt^{-7} -\bq^{-3} \bt^{-9} +\bq^{-5} \bt^{-5} +2 \bq^{-5} \bt^{-7} -3 \bq^{-5} \bt^{-9} +\bq^{-7} \bt^{-7} \\
& \quad -\bq^{-7} \bt^{-11} \\
[\ba^{-3}]\colon\ & -\bq^{6} \bt^{-1} -2 \bq^{4} \bt^{-1} -\bq^{4} \bt^{-3} -\bq^{2} \bt^{-1} -5 \bq^{2} \bt^{-3} +\bq^{2} \bt^{-5} \\
& \quad -6 \bt^{-3} -\bq^{-2} \bt^{-3} -5 \bq^{-2} \bt^{-5} +\bq^{-2} \bt^{-7} -2 \bq^{-4} \bt^{-5} -\bq^{-4} \bt^{-7} \\
& \quad -\bq^{-6} \bt^{-7} \\
[\ba^{-2}]\colon\ & -\bq^{-1} \bt^{-3} +\bq^{-1} \bt^{-5} \\
[\ba^{-1}]\colon\ & -\bt^{-1} +\bt^{-3}
\end{align*}

$(1 + t)(\hat g_{(2)} - \hat g_{(1,1)})$（整数 Laurent、290 項）：
\begin{align*}
[\ba^{-12}]\colon\ & -\bq^{13} \bt^{-9} -\bq \bt^{-9} -\bq^{-1} \bt^{-9} +\bq^{-1} \bt^{-11} +\bq^{-3} \bt^{-11} -\bq^{-3} \bt^{-13} \\
& \quad -\bq^{-5} \bt^{-13} +\bq^{-5} \bt^{-15} +\bq^{-7} \bt^{-15} -\bq^{-7} \bt^{-17} -\bq^{-9} \bt^{-17} +\bq^{-9} \bt^{-19} \\
& \quad +\bq^{-11} \bt^{-19} \\
[\ba^{-11}]\colon\ & -\bq^{14} \bt^{-7} +\bq^{14} \bt^{-9} -\bq^{12} \bt^{-7} -\bq^{10} \bt^{-9} -\bq^{8} \bt^{-9} -\bq^{6} \bt^{-9} \\
& \quad -\bq^{4} \bt^{-11} -\bq^{2} \bt^{-11} -\bt^{-13} -\bq^{-2} \bt^{-13} -\bq^{-4} \bt^{-15} -\bq^{-6} \bt^{-15} \\
& \quad -\bq^{-8} \bt^{-19} -\bq^{-10} \bt^{-19} \\
[\ba^{-10}]\colon\ & \bq^{15} \bt^{-7} -\bq^{13} \bt^{-5} +2 \bq^{13} \bt^{-7} -\bq^{11} \bt^{-5} +2 \bq^{11} \bt^{-9} -\bq^{9} \bt^{-5} \\
& \quad +\bq^{9} \bt^{-9} -\bq^{7} \bt^{-5} -\bq^{7} \bt^{-7} +2 \bq^{7} \bt^{-9} -\bq^{5} \bt^{-7} +\bq^{5} \bt^{-11} \\
& \quad -\bq^{3} \bt^{-5} -\bq^{3} \bt^{-7} +2 \bq^{3} \bt^{-11} -\bq \bt^{-5} +\bq \bt^{-7} -\bq \bt^{-11} \\
& \quad +\bq \bt^{-13} +\bq^{-1} \bt^{-7} -2 \bq^{-1} \bt^{-9} -2 \bq^{-1} \bt^{-11} +3 \bq^{-1} \bt^{-13} -\bq^{-3} \bt^{-9} \\
& \quad +\bq^{-3} \bt^{-13} -\bq^{-5} \bt^{-13} +\bq^{-5} \bt^{-17} -\bq^{-7} \bt^{-15} +\bq^{-7} \bt^{-17} -\bq^{-9} \bt^{-15} \\
& \quad -\bq^{-11} \bt^{-19} -\bq^{-13} \bt^{-19} \\
[\ba^{-9}]\colon\ & 2 \bq^{14} \bt^{-5} -\bq^{14} \bt^{-7} +2 \bq^{12} \bt^{-5} +\bq^{12} \bt^{-7} -\bq^{12} \bt^{-9} +\bq^{10} \bt^{-5} \\
& \quad +2 \bq^{10} \bt^{-7} +\bq^{8} \bt^{-5} +3 \bq^{8} \bt^{-7} +4 \bq^{6} \bt^{-7} +\bq^{6} \bt^{-9} +2 \bq^{4} \bt^{-7} \\
& \quad +4 \bq^{4} \bt^{-9} -\bq^{4} \bt^{-11} +\bq^{2} \bt^{-7} +4 \bq^{2} \bt^{-9} +\bq^{2} \bt^{-11} +3 \bt^{-9} \\
& \quad +4 \bt^{-11} -\bt^{-13} +\bq^{-2} \bt^{-9} +3 \bq^{-2} \bt^{-11} +\bq^{-2} \bt^{-13} +\bq^{-4} \bt^{-11} \\
& \quad +4 \bq^{-4} \bt^{-13} +2 \bq^{-6} \bt^{-13} +2 \bq^{-6} \bt^{-15} +2 \bq^{-8} \bt^{-15} +\bq^{-8} \bt^{-17} +\bq^{-10} \bt^{-17} \\
& \quad +\bq^{-10} \bt^{-19} +\bq^{-12} \bt^{-19} \\
[\ba^{-8}]\colon\ & -\bq^{15} \bt^{-5} +\bq^{13} \bt^{-3} -\bq^{13} \bt^{-5} -\bq^{13} \bt^{-7} +\bq^{11} \bt^{-3} +\bq^{11} \bt^{-5} \\
& \quad -3 \bq^{11} \bt^{-7} +2 \bq^{9} \bt^{-3} +\bq^{9} \bt^{-5} -2 \bq^{9} \bt^{-7} -\bq^{9} \bt^{-9} +\bq^{7} \bt^{-3} \\
& \quad +4 \bq^{7} \bt^{-5} -4 \bq^{7} \bt^{-7} -\bq^{7} \bt^{-9} +6 \bq^{5} \bt^{-5} -\bq^{5} \bt^{-7} -4 \bq^{5} \bt^{-9} \\
& \quad +4 \bq^{3} \bt^{-5} +\bq^{3} \bt^{-7} -4 \bq^{3} \bt^{-9} -\bq^{3} \bt^{-11} +2 \bq \bt^{-5} +4 \bq \bt^{-7} \\
& \quad -\bq \bt^{-9} -4 \bq \bt^{-11} +\bq^{-1} \bt^{-5} +\bq^{-1} \bt^{-7} +3 \bq^{-1} \bt^{-9} -4 \bq^{-1} \bt^{-11} \\
& \quad -\bq^{-1} \bt^{-13} -\bq^{-3} \bt^{-7} +6 \bq^{-3} \bt^{-9} -4 \bq^{-3} \bt^{-13} +3 \bq^{-5} \bt^{-9} -2 \bq^{-5} \bt^{-13} \\
& \quad -\bq^{-5} \bt^{-15} -\bq^{-7} \bt^{-11} +3 \bq^{-7} \bt^{-13} -\bq^{-7} \bt^{-15} -\bq^{-7} \bt^{-17} +2 \bq^{-9} \bt^{-13} \\
& \quad -\bq^{-9} \bt^{-17} +\bq^{-11} \bt^{-17} +\bq^{-13} \bt^{-17} \\
[\ba^{-7}]\colon\ & -\bq^{14} \bt^{-3} -\bq^{12} \bt^{-3} -2 \bq^{12} \bt^{-5} +\bq^{12} \bt^{-7} -2 \bq^{10} \bt^{-3} -2 \bq^{10} \bt^{-5} \\
& \quad -2 \bq^{8} \bt^{-3} -4 \bq^{8} \bt^{-5} -7 \bq^{6} \bt^{-5} -2 \bq^{6} \bt^{-7} +\bq^{6} \bt^{-9} -5 \bq^{4} \bt^{-5} \\
& \quad -6 \bq^{4} \bt^{-7} +\bq^{4} \bt^{-9} -\bq^{2} \bt^{-5} -9 \bq^{2} \bt^{-7} -2 \bq^{2} \bt^{-9} +\bq^{2} \bt^{-11} \\
& \quad +\bt^{-5} -7 \bt^{-7} -6 \bt^{-9} +\bt^{-11} -\bq^{-2} \bt^{-7} -7 \bq^{-2} \bt^{-9} \\
& \quad -3 \bq^{-2} \bt^{-11} +\bq^{-2} \bt^{-13} -3 \bq^{-4} \bt^{-9} -5 \bq^{-4} \bt^{-11} -2 \bq^{-6} \bt^{-11} -4 \bq^{-6} \bt^{-13} \\
& \quad -3 \bq^{-8} \bt^{-13} -\bq^{-8} \bt^{-15} -\bq^{-10} \bt^{-15} -\bq^{-10} \bt^{-17} -\bq^{-12} \bt^{-17} \\
[\ba^{-6}]\colon\ & \bq^{13} \bt^{-5} -\bq^{11} \bt^{-3} +\bq^{11} \bt^{-5} -\bq^{9} \bt^{-1} -\bq^{9} \bt^{-3} +\bq^{9} \bt^{-5} \\
& \quad +\bq^{9} \bt^{-7} -4 \bq^{7} \bt^{-3} +2 \bq^{7} \bt^{-5} +2 \bq^{7} \bt^{-7} -4 \bq^{5} \bt^{-3} -\bq^{5} \bt^{-5} \\
& \quad +5 \bq^{5} \bt^{-7} -7 \bq^{3} \bt^{-5} +5 \bq^{3} \bt^{-7} +2 \bq^{3} \bt^{-9} +\bq \bt^{-3} -9 \bq \bt^{-5} \\
& \quad +3 \bq \bt^{-7} +5 \bq \bt^{-9} -\bq^{-1} \bt^{-5} -4 \bq^{-1} \bt^{-7} +3 \bq^{-1} \bt^{-9} +2 \bq^{-1} \bt^{-11} \\
& \quad +\bq^{-3} \bt^{-5} -5 \bq^{-3} \bt^{-7} -2 \bq^{-3} \bt^{-9} +6 \bq^{-3} \bt^{-11} -\bq^{-5} \bt^{-7} -3 \bq^{-5} \bt^{-9} \\
& \quad +3 \bq^{-5} \bt^{-11} +\bq^{-5} \bt^{-13} -\bq^{-7} \bt^{-9} -\bq^{-7} \bt^{-11} +\bq^{-7} \bt^{-13} +\bq^{-7} \bt^{-15} \\
& \quad -\bq^{-9} \bt^{-13} +\bq^{-9} \bt^{-15} -\bq^{-11} \bt^{-13} \\
[\ba^{-5}]\colon\ & \bq^{12} \bt^{-3} +\bq^{10} \bt^{-1} +\bq^{10} \bt^{-3} +\bq^{8} \bt^{-1} +3 \bq^{8} \bt^{-3} +5 \bq^{6} \bt^{-3} \\
& \quad +2 \bq^{6} \bt^{-5} -\bq^{6} \bt^{-7} +2 \bq^{4} \bt^{-3} +7 \bq^{4} \bt^{-5} -\bq^{4} \bt^{-7} -\bq^{2} \bt^{-3} \\
& \quad +10 \bq^{2} \bt^{-5} +\bq^{2} \bt^{-7} -\bq^{2} \bt^{-9} -\bt^{-3} +4 \bt^{-5} +7 \bt^{-7} \\
& \quad -\bt^{-9} -\bq^{-2} \bt^{-5} +8 \bq^{-2} \bt^{-7} +2 \bq^{-2} \bt^{-9} -\bq^{-2} \bt^{-11} +2 \bq^{-4} \bt^{-7} \\
& \quad +4 \bq^{-4} \bt^{-9} +2 \bq^{-6} \bt^{-9} +2 \bq^{-6} \bt^{-11} +\bq^{-8} \bt^{-11} +\bq^{-8} \bt^{-13} +\bq^{-10} \bt^{-13} \\
[\ba^{-4}]\colon\ & \bq^{7} \bt^{-1} -\bq^{7} \bt^{-5} -\bq^{5} \bt^{-1} +2 \bq^{5} \bt^{-3} -2 \bq^{5} \bt^{-5} -\bq^{3} \bt^{-1} \\
& \quad +5 \bq^{3} \bt^{-3} -3 \bq^{3} \bt^{-5} -\bq^{3} \bt^{-7} +\bq \bt^{-3} +\bq \bt^{-5} -3 \bq \bt^{-7} \\
& \quad -\bq^{-1} \bt^{-3} +5 \bq^{-1} \bt^{-5} -3 \bq^{-1} \bt^{-7} -\bq^{-1} \bt^{-9} +\bq^{-3} \bt^{-7} -2 \bq^{-3} \bt^{-9} \\
& \quad -\bq^{-5} \bt^{-5} +3 \bq^{-5} \bt^{-7} -\bq^{-5} \bt^{-9} -\bq^{-5} \bt^{-11} +\bq^{-7} \bt^{-7} -\bq^{-7} \bt^{-11} \\
[\ba^{-3}]\colon\ & -\bq^{8} \bt^{-1} -\bq^{6} \bt^{-1} -\bq^{6} \bt^{-3} +\bq^{4} \bt^{-1} -4 \bq^{4} \bt^{-3} +\bq^{4} \bt^{-5} \\
& \quad +\bq^{2} \bt^{-1} -3 \bq^{2} \bt^{-3} -\bq^{2} \bt^{-5} +\bt^{-3} -5 \bt^{-5} +\bt^{-7} \\
& \quad +\bq^{-2} \bt^{-3} -3 \bq^{-2} \bt^{-5} -2 \bq^{-4} \bt^{-7} -\bq^{-6} \bt^{-7} \\
[\ba^{-2}]\colon\ & \bq \bt^{-1} -2 \bq \bt^{-3} +\bq \bt^{-5} -\bq^{-1} \bt^{-3} +\bq^{-1} \bt^{-5} \\
[\ba^{-1}]\colon\ & -\bq^{2} \bt^{-1} +\bq^{2} \bt^{-3} -\bt^{-1} +\bt^{-3} \\
[\ba^{0}]\colon\ & \bq \bt^{-1}
\end{align*}

和の商 $X/((q - t)(\nu^2 - t))$（189 項）：
\begin{align*}
[\ba^{-12}]\colon\ & -\bq^{7} \bt^{-3} -\bq^{7} \bt^{-5} -\bq^{7} \bt^{-7} -\bq^{7} \bt^{-9} -\bq^{5} \bt^{-5} -\bq^{5} \bt^{-7} \\
& \quad -\bq^{5} \bt^{-9} -\bq^{3} \bt^{-7} -\bq^{3} \bt^{-9} -\bq \bt^{-9} +\bq^{-3} \bt^{-11} -\bq^{-5} \bt^{-11} \\
& \quad -\bq^{-5} \bt^{-13} -\bq^{-7} \bt^{-13} -\bq^{-7} \bt^{-15} -\bq^{-9} \bt^{-15} -\bq^{-9} \bt^{-17} -\bq^{-11} \bt^{-17} \\
& \quad -\bq^{-11} \bt^{-19} -\bq^{-13} \bt^{-19} -\bq^{-13} \bt^{-21} -\bq^{-15} \bt^{-21} -\bq^{-15} \bt^{-23} \\
[\ba^{-11}]\colon\ & \bq^{8} \bt^{-9} -\bq^{6} \bt^{-3} -\bq^{6} \bt^{-5} -\bq^{6} \bt^{-7} +\bq^{6} \bt^{-9} -2 \bq^{4} \bt^{-5} \\
& \quad -2 \bq^{4} \bt^{-7} -3 \bq^{2} \bt^{-7} -\bq^{2} \bt^{-9} -2 \bt^{-9} +2 \bq^{-2} \bt^{-9} -\bq^{-2} \bt^{-11} \\
& \quad +4 \bq^{-4} \bt^{-11} +4 \bq^{-6} \bt^{-13} +4 \bq^{-8} \bt^{-15} +3 \bq^{-10} \bt^{-17} +2 \bq^{-12} \bt^{-19} +\bq^{-14} \bt^{-21} \\
[\ba^{-10}]\colon\ & \bq^{9} \bt +\bq^{9} \bt^{-1} +\bq^{9} \bt^{-3} +\bq^{9} \bt^{-5} +\bq^{9} \bt^{-7} +2 \bq^{7} \bt^{-1} \\
& \quad +2 \bq^{7} \bt^{-3} +2 \bq^{7} \bt^{-5} +3 \bq^{7} \bt^{-7} +2 \bq^{5} \bt^{-3} +2 \bq^{5} \bt^{-5} +4 \bq^{5} \bt^{-7} \\
& \quad +\bq^{5} \bt^{-9} +2 \bq^{3} \bt^{-5} +5 \bq^{3} \bt^{-7} +2 \bq^{3} \bt^{-9} +6 \bq \bt^{-7} +4 \bq \bt^{-9} \\
& \quad +4 \bq^{-1} \bt^{-7} +5 \bq^{-1} \bt^{-9} +\bq^{-1} \bt^{-11} +2 \bq^{-3} \bt^{-7} +7 \bq^{-3} \bt^{-9} +3 \bq^{-3} \bt^{-11} \\
& \quad +4 \bq^{-5} \bt^{-9} +7 \bq^{-5} \bt^{-11} +\bq^{-5} \bt^{-13} +4 \bq^{-7} \bt^{-11} +6 \bq^{-7} \bt^{-13} +3 \bq^{-9} \bt^{-13} \\
& \quad +4 \bq^{-9} \bt^{-15} +3 \bq^{-11} \bt^{-15} +3 \bq^{-11} \bt^{-17} +2 \bq^{-13} \bt^{-17} +2 \bq^{-13} \bt^{-19} +\bq^{-15} \bt^{-19} \\
& \quad +\bq^{-15} \bt^{-21} \\
[\ba^{-9}]\colon\ & \bq^{8} \bt +\bq^{8} \bt^{-1} +\bq^{8} \bt^{-3} +\bq^{8} \bt^{-5} -\bq^{8} \bt^{-7} +3 \bq^{6} \bt^{-1} \\
& \quad +3 \bq^{6} \bt^{-3} +3 \bq^{6} \bt^{-5} -\bq^{6} \bt^{-7} +5 \bq^{4} \bt^{-3} +5 \bq^{4} \bt^{-5} +8 \bq^{2} \bt^{-5} \\
& \quad +2 \bq^{2} \bt^{-7} -\bq^{2} \bt^{-9} +6 \bt^{-7} -\bt^{-9} -4 \bq^{-2} \bt^{-7} +\bq^{-2} \bt^{-9} \\
& \quad -\bq^{-2} \bt^{-11} -10 \bq^{-4} \bt^{-9} -\bq^{-4} \bt^{-11} -11 \bq^{-6} \bt^{-11} -\bq^{-6} \bt^{-13} -9 \bq^{-8} \bt^{-13} \\
& \quad -6 \bq^{-10} \bt^{-15} -3 \bq^{-12} \bt^{-17} -\bq^{-14} \bt^{-19} \\
[\ba^{-8}]\colon\ & -\bq^{9} \bt^{3} -\bq^{9} \bt -\bq^{9} \bt^{-1} -\bq^{9} \bt^{-3} -\bq^{9} \bt^{-5} -\bq^{7} \bt \\
& \quad -\bq^{7} \bt^{-1} -\bq^{7} \bt^{-3} -2 \bq^{7} \bt^{-5} -\bq^{5} \bt^{-1} -\bq^{5} \bt^{-3} -3 \bq^{5} \bt^{-5} \\
& \quad -\bq^{5} \bt^{-7} -\bq^{3} \bt^{-3} -5 \bq^{3} \bt^{-5} -3 \bq^{3} \bt^{-7} -6 \bq \bt^{-5} -7 \bq \bt^{-7} \\
& \quad -4 \bq^{-1} \bt^{-5} -11 \bq^{-1} \bt^{-7} -2 \bq^{-1} \bt^{-9} -13 \bq^{-3} \bt^{-7} -6 \bq^{-3} \bt^{-9} -3 \bq^{-5} \bt^{-7} \\
& \quad -12 \bq^{-5} \bt^{-9} -2 \bq^{-5} \bt^{-11} -4 \bq^{-7} \bt^{-9} -8 \bq^{-7} \bt^{-11} -3 \bq^{-9} \bt^{-11} -4 \bq^{-9} \bt^{-13} \\
& \quad -2 \bq^{-11} \bt^{-13} -2 \bq^{-11} \bt^{-15} -\bq^{-13} \bt^{-15} -\bq^{-13} \bt^{-17} \\
[\ba^{-7}]\colon\ & -\bq^{8} \bt^{3} -\bq^{8} \bt -\bq^{8} \bt^{-1} -\bq^{8} \bt^{-3} -2 \bq^{6} \bt -2 \bq^{6} \bt^{-1} \\
& \quad -2 \bq^{6} \bt^{-3} -4 \bq^{4} \bt^{-1} -4 \bq^{4} \bt^{-3} -8 \bq^{2} \bt^{-3} -2 \bq^{2} \bt^{-5} +\bq^{2} \bt^{-7} \\
& \quad -7 \bt^{-5} +2 \bt^{-7} +2 \bq^{-2} \bt^{-5} +\bq^{-2} \bt^{-7} +\bq^{-2} \bt^{-9} +7 \bq^{-4} \bt^{-7} \\
& \quad +2 \bq^{-4} \bt^{-9} +9 \bq^{-6} \bt^{-9} +\bq^{-6} \bt^{-11} +6 \bq^{-8} \bt^{-11} +3 \bq^{-10} \bt^{-13} +\bq^{-12} \bt^{-15} \\
[\ba^{-6}]\colon\ & \bq^{3} \bt^{-3} +\bq^{3} \bt^{-5} +\bq \bt^{-3} +3 \bq \bt^{-5} +6 \bq^{-1} \bt^{-5} +\bq^{-1} \bt^{-7} \\
& \quad +3 \bq^{-3} \bt^{-5} +4 \bq^{-3} \bt^{-7} +6 \bq^{-5} \bt^{-7} +\bq^{-5} \bt^{-9} +\bq^{-7} \bt^{-7} +3 \bq^{-7} \bt^{-9} \\
& \quad +\bq^{-9} \bt^{-9} +\bq^{-9} \bt^{-11} \\
[\ba^{-5}]\colon\ & \bq^{4} \bt +\bq^{4} \bt^{-1} +3 \bq^{2} \bt^{-1} +\bq^{2} \bt^{-3} +3 \bt^{-3} -\bt^{-5} \\
& \quad -\bq^{-4} \bt^{-5} -\bq^{-4} \bt^{-7} -2 \bq^{-6} \bt^{-7} -\bq^{-8} \bt^{-9} \\
[\ba^{-4}]\colon\ & -\bq^{-3} \bt^{-5} \\
[\ba^{-3}]\colon\ & -\bq^{-2} \bt^{-3}
\end{align*}

$\hat g_{(1)}$ の商：
\begin{align*}
X_{(1)}/((q - t)(\nu^2 - t)) &= -\ba^{-4} \bq^{-4} \bt^{-4} +\ba^{-5} \bq^{-3} \bt^{-4} +\ba^{-6} \bq^{-4} \bt^{-6}
\end{align*}


\section{被覆の brane--brane と brane--image のセクター}\label{app:cover}
§\ref{sec:split} と §\ref{sec:kn} の根拠。trefoil と $T(2,5)$ の 2 箱について、被覆を brane（$Y$）と image（$Y'$）に分けた生成関数の各セクターを、\cite{KN17} の t 側 $(t^{-1/2} - t^{1/2})\hat f$ で $G$ 展開したもの（$\lambda^\beta$ ごと、$\qr = (q/t)^{1/2}$）。brane--brane は $\hat f_{\ytwo}$, $\hat f_{\yoneone}$（$M \otimes M$ の hatted。一様な規格化の因子を除けば \cite{KN17} の $\hat f$ の鏡像に $R \leftrightarrow R^T$ で一致する、§\ref{sec:split}）、brane--image は $x = W_{\mathrm{adj}} + \kappa - W_{\yone}^2$ を $z_q$ で割ったもの。ブロックの符号の判定は次の表のとおりで、brane--brane は全ブロックでそろい、brane--image は「そろう」「$(\bt^2 - 1)\times$ そろう」「混在」が混じる。
\begin{center}
\footnotesize
\begin{tabular}{llp{9.2cm}}
\toprule
結び目 & セクター & $\lambda^\beta$ ごとの符号 \\
\midrule
$T(2,3)$ & brane--brane $(t^{-1/2} - t^{1/2})\hat f_{\ytwo}$ & $\lambda^{-10}$: そろう、$\lambda^{-8}$: そろう、$\lambda^{-6}$: そろう、$\lambda^{-4}$: そろう、$\lambda^{-2}$: そろう \\
$T(2,3)$ & brane--brane $(t^{-1/2} - t^{1/2})\hat f_{\yoneone}$ & $\lambda^{-10}$: そろう、$\lambda^{-8}$: そろう、$\lambda^{-6}$: そろう、$\lambda^{-4}$: そろう、$\lambda^{-2}$: そろう \\
$T(2,3)$ & brane--image $(t^{-1/2} - t^{1/2})\,x/z_q$ & $\lambda^{-12}$: そろう、$\lambda^{-10}$: $(\bt^2 - 1)\times$ そろう、$\lambda^{-8}$: \textbf{混在}、$\lambda^{-6}$: そろう、$\lambda^{-4}$: そろう、$\lambda^{-2}$: $(\bt^2 - 1)\times$ そろう、$\lambda^{0}$: そろう \\
$T(2,5)$ & brane--brane $(t^{-1/2} - t^{1/2})\hat f_{\ytwo}$ & $\lambda^{-14}$: そろう、$\lambda^{-12}$: そろう、$\lambda^{-10}$: そろう、$\lambda^{-8}$: そろう、$\lambda^{-6}$: そろう \\
$T(2,5)$ & brane--brane $(t^{-1/2} - t^{1/2})\hat f_{\yoneone}$ & $\lambda^{-14}$: そろう、$\lambda^{-12}$: そろう、$\lambda^{-10}$: そろう、$\lambda^{-8}$: そろう、$\lambda^{-6}$: そろう \\
$T(2,5)$ & brane--image $(t^{-1/2} - t^{1/2})\,x/z_q$ & $\lambda^{-20}$: そろう、$\lambda^{-18}$: $(\bt^2 - 1)\times$ そろう、$\lambda^{-16}$: \textbf{混在}、$\lambda^{-14}$: \textbf{混在}、$\lambda^{-12}$: \textbf{混在}、$\lambda^{-10}$: \textbf{混在}、$\lambda^{-8}$: そろう、$\lambda^{-6}$: $(\bt^2 - 1)\times$ そろう、$\lambda^{-4}$: $(\bt^2 - 1)\times$ そろう、$\lambda^{-2}$: $(\bt^2 - 1)\times$ そろう、$\lambda^{0}$: そろう \\
\bottomrule
\end{tabular}
\end{center}

$T(2,3)$、brane--brane $(t^{-1/2} - t^{1/2})\hat f_{\ytwo}$ の $G$ 展開（$\lambda^\beta$ ごと、46 項）
\begin{align*}
[\lambda^{-10}]\colon\ & -G^{2} \qr^{-7} -2 G \qr^{-6} -2 G \qr^{-8} -2 \qr^{-5} -\qr^{-7} -\qr^{-9} \\
[\lambda^{-8}]\colon\ & G^{3} \qr^{-6} +4 G^{2} \qr^{-5} +4 G^{2} \qr^{-7} +6 G \qr^{-4} +9 G \qr^{-6} +5 G \qr^{-8} \\
& \quad +3 \qr^{-3} +7 \qr^{-5} +4 \qr^{-7} +2 \qr^{-9} \\
[\lambda^{-6}]\colon\ & -G^{3} \qr^{-4} -G^{3} \qr^{-6} -3 G^{2} \qr^{-3} -8 G^{2} \qr^{-5} -3 G^{2} \qr^{-7} -3 G \qr^{-2} \\
& \quad -14 G \qr^{-4} -12 G \qr^{-6} -3 G \qr^{-8} -\qr^{-1} -8 \qr^{-3} -9 \qr^{-5} \\
& \quad -5 \qr^{-7} -\qr^{-9} \\
[\lambda^{-4}]\colon\ & G^{3} \qr^{-4} +4 G^{2} \qr^{-3} +4 G^{2} \qr^{-5} +5 G \qr^{-2} +10 G \qr^{-4} +5 G \qr^{-6} \\
& \quad +2 \qr^{-1} +7 \qr^{-3} +5 \qr^{-5} +2 \qr^{-7} \\
[\lambda^{-2}]\colon\ & -G^{2} \qr^{-3} -2 G \qr^{-2} -2 G \qr^{-4} -\qr^{-1} -2 \qr^{-3} -\qr^{-5}
\end{align*}

$T(2,3)$、brane--brane $(t^{-1/2} - t^{1/2})\hat f_{\yoneone}$ の $G$ 展開（$\lambda^\beta$ ごと、27 項）
\begin{align*}
[\lambda^{-10}]\colon\ & -G \qr^{-6} -\qr^{-5} -\qr^{-7} \\
[\lambda^{-8}]\colon\ & G^{2} \qr^{-5} +3 G \qr^{-4} +3 G \qr^{-6} +2 \qr^{-3} +4 \qr^{-5} +2 \qr^{-7} \\
[\lambda^{-6}]\colon\ & -G^{2} \qr^{-3} -G^{2} \qr^{-5} -2 G \qr^{-2} -6 G \qr^{-4} -2 G \qr^{-6} -\qr^{-1} \\
& \quad -5 \qr^{-3} -5 \qr^{-5} -\qr^{-7} \\
[\lambda^{-4}]\colon\ & G^{2} \qr^{-3} +3 G \qr^{-2} +3 G \qr^{-4} +2 \qr^{-1} +4 \qr^{-3} +2 \qr^{-5} \\
[\lambda^{-2}]\colon\ & -G \qr^{-2} -\qr^{-1} -\qr^{-3}
\end{align*}

$T(2,3)$、brane--image $(t^{-1/2} - t^{1/2})\,x/z_q$ の $G$ 展開（$\lambda^\beta$ ごと、29 項）
\begin{align*}
[\lambda^{-12}]\colon\ & -\qr^{-9} \\
[\lambda^{-10}]\colon\ & -G \qr^{-6} +G \qr^{-8} -2 \qr^{-5} +2 \qr^{-9} \\
[\lambda^{-8}]\colon\ & G^{2} \qr^{-5} +4 G \qr^{-4} +3 G \qr^{-6} -G \qr^{-8} +3 \qr^{-3} +7 \qr^{-5} \\
& \quad -\qr^{-9} \\
[\lambda^{-6}]\colon\ & -G^{2} \qr^{-3} -G^{2} \qr^{-5} -2 G \qr^{-2} -8 G \qr^{-4} -2 G \qr^{-6} -\qr^{-1} \\
& \quad -7 \qr^{-3} -8 \qr^{-5} \\
[\lambda^{-4}]\colon\ & G^{2} \qr^{-3} +2 G \qr^{-2} +4 G \qr^{-4} +\qr^{-1} +5 \qr^{-3} +3 \qr^{-5} \\
[\lambda^{-2}]\colon\ & \qr^{-1} -\qr^{-3} \\
[\lambda^{0}]\colon\ & -\qr^{-1}
\end{align*}

$T(2,3)$、同じものを $\lambda^\beta$ ごとに因数分解した形（$\qr^2 - 1 = \bt^2 - 1$）
\begin{align*}
[\lambda^{-12}]\colon\ & \qr^{-9}\,\bigl(-1\bigr) \\
[\lambda^{-10}]\colon\ & \qr^{-9}\,\bigl(- (\qr^2 - 1)(G \qr + 2 \qr^{2} + 2)\bigr) \\
[\lambda^{-8}]\colon\ & \text{既約（上の形のまま）} \\
[\lambda^{-6}]\colon\ & \text{既約（上の形のまま）} \\
[\lambda^{-4}]\colon\ & \text{既約（上の形のまま）} \\
[\lambda^{-2}]\colon\ & \qr^{-3}\,(\qr^2 - 1) \\
[\lambda^{0}]\colon\ & \qr^{-1}\,\bigl(-1\bigr)
\end{align*}

$T(2,5)$、brane--brane $(t^{-1/2} - t^{1/2})\hat f_{\ytwo}$ の $G$ 展開（$\lambda^\beta$ ごと、172 項）
\begin{align*}
[\lambda^{-14}]\colon\ & -G^{6} \qr^{-11} -6 G^{5} \qr^{-10} -6 G^{5} \qr^{-12} -16 G^{4} \qr^{-9} -26 G^{4} \qr^{-11} -15 G^{4} \qr^{-13} \\
& \quad -24 G^{3} \qr^{-8} -49 G^{3} \qr^{-10} -44 G^{3} \qr^{-12} -20 G^{3} \qr^{-14} -22 G^{2} \qr^{-7} -49 G^{2} \qr^{-9} \\
& \quad -54 G^{2} \qr^{-11} -36 G^{2} \qr^{-13} -15 G^{2} \qr^{-15} -12 G \qr^{-6} -28 G \qr^{-8} -30 G \qr^{-10} \\
& \quad -25 G \qr^{-12} -14 G \qr^{-14} -6 G \qr^{-16} -3 \qr^{-5} -7 \qr^{-7} -8 \qr^{-9} \\
& \quad -5 \qr^{-11} -4 \qr^{-13} -2 \qr^{-15} -\qr^{-17} \\
[\lambda^{-12}]\colon\ & G^{7} \qr^{-10} +8 G^{6} \qr^{-9} +8 G^{6} \qr^{-11} +28 G^{5} \qr^{-8} +50 G^{5} \qr^{-10} +27 G^{5} \qr^{-12} \\
& \quad +56 G^{4} \qr^{-7} +133 G^{4} \qr^{-9} +126 G^{4} \qr^{-11} +50 G^{4} \qr^{-13} +70 G^{3} \qr^{-6} +195 G^{3} \qr^{-8} \\
& \quad +240 G^{3} \qr^{-10} +164 G^{3} \qr^{-12} +55 G^{3} \qr^{-14} +55 G^{2} \qr^{-5} +170 G^{2} \qr^{-7} +237 G^{2} \qr^{-9} \\
& \quad +205 G^{2} \qr^{-11} +116 G^{2} \qr^{-13} +36 G^{2} \qr^{-15} +25 G \qr^{-4} +84 G \qr^{-6} +127 G \qr^{-8} \\
& \quad +117 G \qr^{-10} +82 G \qr^{-12} +42 G \qr^{-14} +13 G \qr^{-16} +5 \qr^{-3} +18 \qr^{-5} \\
& \quad +30 \qr^{-7} +28 \qr^{-9} +19 \qr^{-11} +12 \qr^{-13} +6 \qr^{-15} +2 \qr^{-17} \\
[\lambda^{-10}]\colon\ & -G^{7} \qr^{-8} -G^{7} \qr^{-10} -7 G^{6} \qr^{-7} -16 G^{6} \qr^{-9} -7 G^{6} \qr^{-11} -22 G^{5} \qr^{-6} \\
& \quad -72 G^{5} \qr^{-8} -71 G^{5} \qr^{-10} -21 G^{5} \qr^{-12} -40 G^{4} \qr^{-5} -163 G^{4} \qr^{-7} -228 G^{4} \qr^{-9} \\
& \quad -150 G^{4} \qr^{-11} -35 G^{4} \qr^{-13} -45 G^{3} \qr^{-4} -216 G^{3} \qr^{-6} -368 G^{3} \qr^{-8} -335 G^{3} \qr^{-10} \\
& \quad -175 G^{3} \qr^{-12} -35 G^{3} \qr^{-14} -31 G^{2} \qr^{-3} -172 G^{2} \qr^{-5} -334 G^{2} \qr^{-7} -360 G^{2} \qr^{-9} \\
& \quad -252 G^{2} \qr^{-11} -116 G^{2} \qr^{-13} -21 G^{2} \qr^{-15} -12 G \qr^{-2} -77 G \qr^{-4} -166 G \qr^{-6} \\
& \quad -196 G \qr^{-8} -158 G \qr^{-10} -93 G \qr^{-12} -41 G \qr^{-14} -7 G \qr^{-16} -2 \qr^{-1} \\
& \quad -15 \qr^{-3} -35 \qr^{-5} -46 \qr^{-7} -37 \qr^{-9} -25 \qr^{-11} -13 \qr^{-13} \\
& \quad -6 \qr^{-15} -\qr^{-17} \\
[\lambda^{-8}]\colon\ & G^{7} \qr^{-8} +8 G^{6} \qr^{-7} +8 G^{6} \qr^{-9} +28 G^{5} \qr^{-6} +50 G^{5} \qr^{-8} +27 G^{5} \qr^{-10} \\
& \quad +56 G^{4} \qr^{-5} +133 G^{4} \qr^{-7} +126 G^{4} \qr^{-9} +50 G^{4} \qr^{-11} +69 G^{3} \qr^{-4} +195 G^{3} \qr^{-6} \\
& \quad +241 G^{3} \qr^{-8} +164 G^{3} \qr^{-10} +55 G^{3} \qr^{-12} +52 G^{2} \qr^{-3} +166 G^{2} \qr^{-5} +241 G^{2} \qr^{-7} \\
& \quad +208 G^{2} \qr^{-9} +116 G^{2} \qr^{-11} +36 G^{2} \qr^{-13} +22 G \qr^{-2} +77 G \qr^{-4} +127 G \qr^{-6} \\
& \quad +124 G \qr^{-8} +85 G \qr^{-10} +42 G \qr^{-12} +13 G \qr^{-14} +4 \qr^{-1} +15 \qr^{-3} \\
& \quad +28 \qr^{-5} +30 \qr^{-7} +22 \qr^{-9} +13 \qr^{-11} +6 \qr^{-13} +2 \qr^{-15} \\
[\lambda^{-6}]\colon\ & -G^{6} \qr^{-7} -6 G^{5} \qr^{-6} -6 G^{5} \qr^{-8} -16 G^{4} \qr^{-5} -26 G^{4} \qr^{-7} -15 G^{4} \qr^{-9} \\
& \quad -24 G^{3} \qr^{-4} -49 G^{3} \qr^{-6} -44 G^{3} \qr^{-8} -20 G^{3} \qr^{-10} -21 G^{2} \qr^{-3} -49 G^{2} \qr^{-5} \\
& \quad -55 G^{2} \qr^{-7} -36 G^{2} \qr^{-9} -15 G^{2} \qr^{-11} -10 G \qr^{-2} -25 G \qr^{-4} -33 G \qr^{-6} \\
& \quad -27 G \qr^{-8} -14 G \qr^{-10} -6 G \qr^{-12} -2 \qr^{-1} -5 \qr^{-3} -8 \qr^{-5} \\
& \quad -7 \qr^{-7} -5 \qr^{-9} -2 \qr^{-11} -\qr^{-13}
\end{align*}

$T(2,5)$、brane--brane $(t^{-1/2} - t^{1/2})\hat f_{\yoneone}$ の $G$ 展開（$\lambda^\beta$ ごと、133 項）
\begin{align*}
[\lambda^{-14}]\colon\ & -G^{5} \qr^{-10} -5 G^{4} \qr^{-9} -5 G^{4} \qr^{-11} -11 G^{3} \qr^{-8} -17 G^{3} \qr^{-10} -10 G^{3} \qr^{-12} \\
& \quad -13 G^{2} \qr^{-7} -25 G^{2} \qr^{-9} -21 G^{2} \qr^{-11} -10 G^{2} \qr^{-13} -8 G \qr^{-6} -18 G \qr^{-8} \\
& \quad -18 G \qr^{-10} -11 G \qr^{-12} -5 G \qr^{-14} -2 \qr^{-5} -6 \qr^{-7} -5 \qr^{-9} \\
& \quad -4 \qr^{-11} -2 \qr^{-13} -\qr^{-15} \\
[\lambda^{-12}]\colon\ & G^{6} \qr^{-9} +7 G^{5} \qr^{-8} +7 G^{5} \qr^{-10} +21 G^{4} \qr^{-7} +37 G^{4} \qr^{-9} +20 G^{4} \qr^{-11} \\
& \quad +35 G^{3} \qr^{-6} +81 G^{3} \qr^{-8} +75 G^{3} \qr^{-10} +30 G^{3} \qr^{-12} +34 G^{2} \qr^{-5} +94 G^{2} \qr^{-7} \\
& \quad +110 G^{2} \qr^{-9} +73 G^{2} \qr^{-11} +25 G^{2} \qr^{-13} +18 G \qr^{-4} +57 G \qr^{-6} +78 G \qr^{-8} \\
& \quad +62 G \qr^{-10} +34 G \qr^{-12} +11 G \qr^{-14} +4 \qr^{-3} +14 \qr^{-5} +23 \qr^{-7} \\
& \quad +19 \qr^{-9} +12 \qr^{-11} +6 \qr^{-13} +2 \qr^{-15} \\
[\lambda^{-10}]\colon\ & -G^{6} \qr^{-7} -G^{6} \qr^{-9} -6 G^{5} \qr^{-6} -14 G^{5} \qr^{-8} -6 G^{5} \qr^{-10} -16 G^{4} \qr^{-5} \\
& \quad -53 G^{4} \qr^{-7} -52 G^{4} \qr^{-9} -15 G^{4} \qr^{-11} -24 G^{3} \qr^{-4} -99 G^{3} \qr^{-6} -135 G^{3} \qr^{-8} \\
& \quad -88 G^{3} \qr^{-10} -20 G^{3} \qr^{-12} -21 G^{2} \qr^{-3} -102 G^{2} \qr^{-5} -171 G^{2} \qr^{-7} -148 G^{2} \qr^{-9} \\
& \quad -77 G^{2} \qr^{-11} -15 G^{2} \qr^{-13} -10 G \qr^{-2} -55 G \qr^{-4} -109 G \qr^{-6} -113 G \qr^{-8} \\
& \quad -73 G \qr^{-10} -34 G \qr^{-12} -6 G \qr^{-14} -2 \qr^{-1} -12 \qr^{-3} -28 \qr^{-5} \\
& \quad -33 \qr^{-7} -25 \qr^{-9} -13 \qr^{-11} -6 \qr^{-13} -\qr^{-15} \\
[\lambda^{-8}]\colon\ & G^{6} \qr^{-7} +7 G^{5} \qr^{-6} +7 G^{5} \qr^{-8} +21 G^{4} \qr^{-5} +37 G^{4} \qr^{-7} +20 G^{4} \qr^{-9} \\
& \quad +35 G^{3} \qr^{-4} +81 G^{3} \qr^{-6} +75 G^{3} \qr^{-8} +30 G^{3} \qr^{-10} +34 G^{2} \qr^{-3} +93 G^{2} \qr^{-5} \\
& \quad +111 G^{2} \qr^{-7} +73 G^{2} \qr^{-9} +25 G^{2} \qr^{-11} +18 G \qr^{-2} +54 G \qr^{-4} +79 G \qr^{-6} \\
& \quad +64 G \qr^{-8} +34 G \qr^{-10} +11 G \qr^{-12} +4 \qr^{-1} +12 \qr^{-3} +22 \qr^{-5} \\
& \quad +21 \qr^{-7} +13 \qr^{-9} +6 \qr^{-11} +2 \qr^{-13} \\
[\lambda^{-6}]\colon\ & -G^{5} \qr^{-6} -5 G^{4} \qr^{-5} -5 G^{4} \qr^{-7} -11 G^{3} \qr^{-4} -17 G^{3} \qr^{-6} -10 G^{3} \qr^{-8} \\
& \quad -13 G^{2} \qr^{-3} -25 G^{2} \qr^{-5} -21 G^{2} \qr^{-7} -10 G^{2} \qr^{-9} -8 G \qr^{-2} -17 G \qr^{-4} \\
& \quad -19 G \qr^{-6} -11 G \qr^{-8} -5 G \qr^{-10} -2 \qr^{-1} -4 \qr^{-3} -6 \qr^{-5} \\
& \quad -5 \qr^{-7} -2 \qr^{-9} -\qr^{-11}
\end{align*}

$T(2,5)$、brane--image $(t^{-1/2} - t^{1/2})\,x/z_q$ の $G$ 展開（$\lambda^\beta$ ごと、93 項）
\begin{align*}
[\lambda^{-20}]\colon\ & -\qr^{-17} \\
[\lambda^{-18}]\colon\ & -G \qr^{-14} +G \qr^{-16} -\qr^{-13} -\qr^{-15} +2 \qr^{-17} \\
[\lambda^{-16}]\colon\ & -G^{2} \qr^{-11} +G^{2} \qr^{-13} -2 G \qr^{-10} -G \qr^{-12} +4 G \qr^{-14} -G \qr^{-16} \\
& \quad -\qr^{-9} -2 \qr^{-11} +2 \qr^{-13} +2 \qr^{-15} -\qr^{-17} \\
[\lambda^{-14}]\colon\ & -G^{3} \qr^{-8} +G^{3} \qr^{-10} -4 G^{2} \qr^{-7} +5 G^{2} \qr^{-11} -G^{2} \qr^{-13} -5 G \qr^{-6} \\
& \quad -5 G \qr^{-8} +8 G \qr^{-10} +5 G \qr^{-12} -3 G \qr^{-14} -2 \qr^{-5} -4 \qr^{-7} \\
& \quad +\qr^{-9} +7 \qr^{-11} -\qr^{-13} -\qr^{-15} \\
[\lambda^{-12}]\colon\ & G^{4} \qr^{-7} +6 G^{3} \qr^{-6} +5 G^{3} \qr^{-8} -G^{3} \qr^{-10} +12 G^{2} \qr^{-5} +23 G^{2} \qr^{-7} \\
& \quad +4 G^{2} \qr^{-9} -4 G^{2} \qr^{-11} +10 G \qr^{-4} +29 G \qr^{-6} +23 G \qr^{-8} -8 G \qr^{-10} \\
& \quad -4 G \qr^{-12} +3 \qr^{-3} +11 \qr^{-5} +15 \qr^{-7} +4 \qr^{-9} -8 \qr^{-11} \\
[\lambda^{-10}]\colon\ & -G^{4} \qr^{-5} -G^{4} \qr^{-7} -4 G^{3} \qr^{-4} -12 G^{3} \qr^{-6} -4 G^{3} \qr^{-8} -6 G^{2} \qr^{-3} \\
& \quad -28 G^{2} \qr^{-5} -32 G^{2} \qr^{-7} -4 G^{2} \qr^{-9} -4 G \qr^{-2} -24 G \qr^{-4} -44 G \qr^{-6} \\
& \quad -30 G \qr^{-8} +2 G \qr^{-10} -\qr^{-1} -7 \qr^{-3} -17 \qr^{-5} -19 \qr^{-7} \\
& \quad -7 \qr^{-9} +3 \qr^{-11} \\
[\lambda^{-8}]\colon\ & G^{4} \qr^{-5} +4 G^{3} \qr^{-4} +6 G^{3} \qr^{-6} +6 G^{2} \qr^{-3} +16 G^{2} \qr^{-5} +13 G^{2} \qr^{-7} \\
& \quad +4 G \qr^{-2} +14 G \qr^{-4} +20 G \qr^{-6} +12 G \qr^{-8} +\qr^{-1} +4 \qr^{-3} \\
& \quad +8 \qr^{-5} +9 \qr^{-7} +3 \qr^{-9} \\
[\lambda^{-6}]\colon\ & \qr^{-5} -\qr^{-7} \\
[\lambda^{-4}]\colon\ & \qr^{-3} -\qr^{-5} \\
[\lambda^{-2}]\colon\ & \qr^{-1} -\qr^{-3} \\
[\lambda^{0}]\colon\ & -\qr^{-1}
\end{align*}

$T(2,5)$、同じものを $\lambda^\beta$ ごとに因数分解した形（$\qr^2 - 1 = \bt^2 - 1$）
\begin{align*}
[\lambda^{-20}]\colon\ & \qr^{-17}\,\bigl(-1\bigr) \\
[\lambda^{-18}]\colon\ & \qr^{-17}\,\bigl(- (\qr^2 - 1)(G \qr + \qr^{2} + 2)\bigr) \\
[\lambda^{-16}]\colon\ & \qr^{-17}\,\bigl(- (\qr^2 - 1)(G \qr + \qr^{2} + 1)(G \qr^{3} + \qr^{4} + 2 \qr^{2} - 1)\bigr) \\
[\lambda^{-14}]\colon\ & \qr^{-15}\,\bigl(- (\qr^2 - 1) \\
& \quad (G^{3} \qr^{5} + 4 G^{2} \qr^{6} + 4 G^{2} \qr^{4} - G^{2} \qr^{2} \\
& \quad  + 5 G \qr^{7} + 10 G \qr^{5} + 2 G \qr^{3} - 3 G \qr + 2 \qr^{8} + 6 \qr^{6} \\
& \quad  + 5 \qr^{4} - 2 \qr^{2} - 1)\bigr) \\
[\lambda^{-12}]\colon\ & \text{既約（上の形のまま）} \\
[\lambda^{-10}]\colon\ & \text{既約（上の形のまま）} \\
[\lambda^{-8}]\colon\ & \text{既約（上の形のまま）} \\
[\lambda^{-6}]\colon\ & \qr^{-7}\,(\qr^2 - 1) \\
[\lambda^{-4}]\colon\ & \qr^{-5}\,(\qr^2 - 1) \\
[\lambda^{-2}]\colon\ & \qr^{-3}\,(\qr^2 - 1) \\
[\lambda^{0}]\colon\ & \qr^{-1}\,\bigl(-1\bigr)
\end{align*}


\section{unrefined の確認}\label{app:unref}
Mariño の trefoil の $\hat g$（標準 framing）を自前のコードで計算した値。$\hat g_{\yoneone}$ は \cite{Marino09} (4.8) の印刷値と全係数一致し、$\hat g_{\ytwo}$ は $z^9$ の項以外が一致する。
$\hat g_{\ytwo}$（$z = t_M - t_M^{-1}$ の冪ごと、$\lambda$ は Mariño の $\nu$）：
\begin{align*}
[z^{0}]\colon\ & -21 \lambda^{11} +79 \lambda^{9} -111 \lambda^{7} +69 \lambda^{5} -16 \lambda^{3} \\
[z]\colon\ & 21 \lambda^{12} -63 \lambda^{10} +63 \lambda^{8} -21 \lambda^{6} \\
[z^{2}]\colon\ & -70 \lambda^{11} +251 \lambda^{9} -307 \lambda^{7} +146 \lambda^{5} -20 \lambda^{3} \\
[z^{3}]\colon\ & 70 \lambda^{12} -231 \lambda^{10} +231 \lambda^{8} -70 \lambda^{6} \\
[z^{4}]\colon\ & -84 \lambda^{11} +330 \lambda^{9} -366 \lambda^{7} +128 \lambda^{5} -8 \lambda^{3} \\
[z^{5}]\colon\ & 84 \lambda^{12} -322 \lambda^{10} +322 \lambda^{8} -84 \lambda^{6} \\
[z^{6}]\colon\ & -45 \lambda^{11} +220 \lambda^{9} -230 \lambda^{7} +56 \lambda^{5} -\lambda^{3} \\
[z^{7}]\colon\ & 45 \lambda^{12} -219 \lambda^{10} +219 \lambda^{8} -45 \lambda^{6} \\
[z^{8}]\colon\ & -11 \lambda^{11} +78 \lambda^{9} -79 \lambda^{7} +12 \lambda^{5} \\
[z^{9}]\colon\ & 11 \lambda^{12} -78 \lambda^{10} +78 \lambda^{8} -11 \lambda^{6} \\
[z^{10}]\colon\ & -\lambda^{11} +14 \lambda^{9} -14 \lambda^{7} +\lambda^{5} \\
[z^{11}]\colon\ & \lambda^{12} -14 \lambda^{10} +14 \lambda^{8} -\lambda^{6} \\
[z^{12}]\colon\ & \lambda^{9} -\lambda^{7} \\
[z^{13}]\colon\ & -\lambda^{10} +\lambda^{8}
\end{align*}

$\hat g_{\yoneone}$（$z = t_M - t_M^{-1}$ の冪ごと、$\lambda$ は Mariño の $\nu$）：
\begin{align*}
[z^{0}]\colon\ & -15 \lambda^{11} +53 \lambda^{9} -69 \lambda^{7} +39 \lambda^{5} -8 \lambda^{3} \\
[z]\colon\ & 15 \lambda^{12} -45 \lambda^{10} +45 \lambda^{8} -15 \lambda^{6} \\
[z^{2}]\colon\ & -35 \lambda^{11} +126 \lambda^{9} -146 \lambda^{7} +61 \lambda^{5} -6 \lambda^{3} \\
[z^{3}]\colon\ & 35 \lambda^{12} -120 \lambda^{10} +120 \lambda^{8} -35 \lambda^{6} \\
[z^{4}]\colon\ & -28 \lambda^{11} +120 \lambda^{9} -128 \lambda^{7} +37 \lambda^{5} -\lambda^{3} \\
[z^{5}]\colon\ & 28 \lambda^{12} -119 \lambda^{10} +119 \lambda^{8} -28 \lambda^{6} \\
[z^{6}]\colon\ & -9 \lambda^{11} +55 \lambda^{9} -56 \lambda^{7} +10 \lambda^{5} \\
[z^{7}]\colon\ & 9 \lambda^{12} -55 \lambda^{10} +55 \lambda^{8} -9 \lambda^{6} \\
[z^{8}]\colon\ & -\lambda^{11} +12 \lambda^{9} -12 \lambda^{7} +\lambda^{5} \\
[z^{9}]\colon\ & \lambda^{12} -12 \lambda^{10} +12 \lambda^{8} -\lambda^{6} \\
[z^{10}]\colon\ & \lambda^{9} -\lambda^{7} \\
[z^{11}]\colon\ & -\lambda^{10} +\lambda^{8}
\end{align*}



% 書誌は refs.bib（biblatex、alphabetic 様式）。書誌は arXiv API（2026-10-20）で確認した。掲載誌は著者が arXiv に記載したもの。無いものは「掲載誌未確認」。
\printbibliography[title={参考文献}]
\noindent Mimachi（双対 Cauchy 恒等式）、Rains・Rains--Warnaar（virtual Koornwinder）、van Diejen--Sahi（Koornwinder の評価式）の書誌は未確認で、本文では名前だけを挙げた。

\end{document}
